\chapter{腿式里程计}
\label{ch:leg}

% ════════════════════════════════════════════════════════════════
%  本章：吸收 SLAM Handbook ch12「Leg Odometry for SLAM」(Camurri & Mattamala)
%  + Bloesch 等 kinematic-inertial 滤波 + Hartley 等接触辅助 InEKF / 接触预积分
%  + Wisth-Camurri-Fallon 速度预积分(VILENS) + Camurri 接触检测综述。
%  定位：P2 末尾「新型传感器与平台里程计」章群（雷达 ch:radar / 事件 ch:event 为兄弟章）。
%  全书三条已有线的汇流点：InEKF(paper:inekf) + SE_2(3)(sec:se23) + 预积分(ch:imu) + 足式控制(sec:ctrl-legged)。
%  自包含、右扰动主线；接触/速度预积分写成 IMU 预积分(eq:imu-preint-model)的同构平行实例。
%  教学层遵循 v5.0 算法工程教学规范。
% ════════════════════════════════════════════════════════════════

\begin{quote}\small\itshape
本章是“感知与 SLAM”部末尾\textbf{新型传感器与平台里程计}章群的收尾一章。前面几章以相机、激光、IMU 为主线，建立了 SLAM 的完整骨架——观测模型、预积分、因子图、流形上的滤波与优化；本章群把同一套数学迁移到三类新兴的感知与平台：毫米波雷达（\cref{ch:radar}，电磁波测距 + 径向速度全天候）、事件相机（\cref{ch:event}，异步亮度变化、超高动态范围）、以及足式机器人（本章，把腿的运动学读成里程计）。它们都不另起炉灶，而是把已学的语言换一种传感器再讲一遍。把腿式里程计放在最后，是因为它\emph{与全书数学衔接最紧}：它正是 \cref{paper:inekf} 的不变卡尔曼滤波、\cref{sec:se23} 的扩展位姿群 $\mathrm{SE}_2(3)$、与 \cref{ch:imu} 的预积分三条线的\textbf{汇流点}，并与 \cref{sec:ctrl-legged} 的足式控制互为表里——读完它，P2 的感知与 P4 的控制在足式平台上接成一个闭环。
\end{quote}

\begin{practice}[前置自测]
开始之前，先试答下列问题。若已能流畅作答，可只速读本章的对照表与速查卡；若卡住，正是本章要为你解决的。
\begin{enumerate}
  \item 一台轮式机器人靠“数轮子转了多少圈”估计走了多远，这叫轮式里程计。一台四足机器人没有轮子，它的“里程计”应该数什么？关键的几何假设是什么？
  \item 给定一条腿所有关节的角度 $\mathbf{q}$，你能算出脚相对机身的位姿吗？这个映射叫什么？它和机械臂的正运动学是一回事吗？
  \item 站立中的脚“踩死”在地上不动。由此能写出机身速度与关节角速度之间的什么关系？为什么这条关系\emph{只}在脚不打滑时成立？
  \item IMU 预积分把两关键帧间上千个惯性采样压成一条相对运动测量（\cref{ch:imu}）。腿的接触/速度测量也是高频的，能不能照搬同一套“预积分”把它压成一条腿式因子？两者结构一样吗？
  \item 怎么知道哪条腿此刻在“站立”（可信）、哪条在“摆动”（不可信）？如果机器人没有专门的接触传感器，只有关节力矩，还能判断吗？
  \item \cref{paper:inekf} 说不变卡尔曼滤波（InEKF）的一个标志性应用是“足式机器人接触辅助惯导”。把脚的位置加进 $\mathrm{SE}_2(3)$ 状态后，为什么误差传播仍能与轨迹无关？
\end{enumerate}
\end{practice}

\section*{本章目标}
学完本章，你应当能够：
\begin{enumerate}
  \item \textbf{建立}足式机器人的参考系与状态，并\textbf{自包含地写出}单腿正运动学 $\mathbf{fk}(\mathbf{q})$ 与雅可比 $\mathbf{J}(\mathbf{q})$——它们与 \cref{sec:ctrl-legged} 足式控制用的是\emph{同一套} $\mathbf{fk}/\mathbf{J}$，估计侧用来推位姿、控制侧用来算力矩；
  \item \textbf{推导}两类腿式里程计测量：基于正运动学的\textbf{相对位姿}（接触系静止假设 $\Rightarrow$ $\mathbf{T}_{b'b}=\mathbf{fk}(\mathbf{q}')^{-1}\mathbf{fk}(\mathbf{q})$）与基于微分运动学的\textbf{机身速度}（接触点零速度约束 $\Rightarrow$ $\mathbf{v}_b^w=-\boldsymbol{\omega}_b^w\times\mathbf{f}_p(\mathbf{q})-\mathbf{J}_v(\mathbf{q})\dot{\mathbf{q}}$）；
  \item \textbf{掌握}接触估计：摩擦圆锥与压力中心的不打滑条件、为何“阈值化法向力 $f_z$”是主流，以及无接触传感器时如何从关节力矩反演地面反力；
  \item \textbf{把腿式因子写成预积分形式}：编码器噪声的一阶传播，正运动学因子、\textbf{接触预积分因子}、\textbf{速度预积分因子}的残差与协方差，并看出它们是 \cref{ch:imu} IMU 预积分的\emph{同构平行实例}——同样把高频测量积成与初值无关的相对增量、同样进因子图、同样用信息矩阵 $\boldsymbol{\Lambda}$ 加权；
  \item \textbf{兑现 \cref{paper:inekf} 的前向许诺}：把接触点位置增广进 $\mathrm{SE}_2(3)$（\cref{sec:se23}），构造接触辅助不变 EKF，解释接触约束为何保持群仿射结构、从而保住 InEKF 的一致性优势；
  \item \textbf{设计}腿-惯-视-激多模态融合（松/紧耦合），把它看作 \cref{ch:vio} 与 \cref{ch:lidar_slam} 因子图的自然推广，并理解腿变形、地面沉陷、打滑等开放问题为何让“刚性接触”假设失效。
\end{enumerate}

\subsection*{本章知识导航}
本章是一条“\emph{从关节角到机身位姿/速度}”的主线：先用\textbf{运动学}把关节读数变成单腿测量，再借\textbf{接触估计}挑出可信的站立腿，最后把这些测量送进\textbf{因子图（预积分）}或\textbf{不变滤波（InEKF）}两类估计器，并与外感受传感器\textbf{融合}。结构如下：

\begin{center}
\small
\begin{tikzpicture}[
  node distance=4mm and 7mm, >=Stealth,
  blk/.style={draw, rounded corners, align=center, font=\small, inner sep=3pt, fill=main!6, draw=main!60!black},
  grp/.style={draw, rounded corners, align=center, font=\small, inner sep=3pt, fill=second!6, draw=second!60!black},
  bk/.style={draw, rounded corners, align=center, font=\small, inner sep=3pt, fill=third!8, draw=third!60!black},
  ln/.style={->, thick, gray!60!black}]
\node[blk] (fk) {腿运动学\\ $\mathbf{fk}(\mathbf{q}),\mathbf{J}(\mathbf{q})$};
\node[blk, right=of fk] (motion) {运动估计\\ 相对位姿 / 速度};
\node[blk, right=of motion] (contact) {接触估计\\ 摩擦锥 / 阈值 $f_z$};
\node[grp, below=8mm of fk] (fkf) {正运动学因子\\ + 接触预积分};
\node[grp, right=of fkf] (velf) {速度预积分因子\\（点足四足）};
\node[grp, right=of velf] (inekf) {接触辅助 InEKF\\ $\mathrm{SE}_2(3)$ + 接触点};
\node[bk, below=8mm of fkf] (fuse) {多传感融合\\ 腿-惯-视-激};
\node[bk, right=of fuse] (loose) {松耦合\\（择优）};
\node[bk, right=of loose] (tight) {紧耦合\\（同一因子图）};
\draw[ln] (fk)--(motion); \draw[ln] (motion)--(contact);
\draw[ln] (motion.south) -- ++(0,-4.5mm) -| (fkf.north);
\draw[ln] (contact.south) -- ++(0,-4.5mm) -| (inekf.north);
\draw[ln] (fkf)--(velf); \draw[ln] (velf)--(inekf);
\draw[ln] (fkf.south) -- (fuse.north);
\draw[ln] (fuse)--(loose); \draw[ln] (loose)--(tight);
\end{tikzpicture}
\end{center}

\noindent\textbf{推荐路径（骨干 only）}：先读 \cref{sec:leg-why}（为什么、两子问题）$\to$ \cref{sec:leg-kinematics}（参考系、状态、$\mathbf{fk}/\mathbf{J}$）$\to$ \cref{sec:leg-motion}（相对位姿与速度，两条核心测量）$\to$ \cref{sec:leg-contact}（接触估计，阈值化 $f_z$）$\to$ \cref{sec:leg-graph}（因子图：正运动学/接触/速度预积分），约 4--5 小时即可建立可实现的完整链条。\cref{sec:leg-inekf}（接触辅助 InEKF）承接 \cref{paper:inekf}/\cref{sec:se23}，理论地基已在书中，可较快读完；\cref{sec:leg-fusion}（多传感融合）与 \cref{sec:leg-open}（开放问题与前沿）面向系统与研究，可选读。

\subsection*{前置知识桥接}
本章把全书几条主线“拼装”成足式平台的里程计，请先激活以下前置：
\begin{itemize}
  \item \textbf{李群与右扰动}（\cref{ch:lie}）：$\mathbf{R}\in\SO(3)$、$\SE(3)$ 齐次变换、$\mathrm{Exp}/\mathrm{Log}$、\textbf{右扰动} $\mathbf{R}\,\mathrm{Exp}(\delta\boldsymbol{\phi})$ 与右雅可比 $\mathbf{J}_r$。腿的正运动学输出 $\SE(3)$ 位姿，相对位姿与因子残差都建在这套代数上。
  \item \textbf{扩展位姿群 $\mathrm{SE}_2(3)$}（\cref{sec:se23}，矩阵 $[\mathbf{R},\mathbf{r},\mathbf{v}]$）：本章接触辅助 InEKF 的状态就活在 $\mathrm{SE}_2(3)$（再增广接触点位置）上，直接复用 $\boldsymbol{\xi}=[\boldsymbol{\rho};\boldsymbol{\nu};\boldsymbol{\phi}]$ 与右扰动约定，零记号摩擦。
  \item \textbf{IMU 模型与预积分}（\cref{ch:imu}）：预积分把高频测量积成与初态无关的相对增量、放进因子图、用一阶传播算协方差（\cref{eq:imu-preint-model,eq:cov-recursion}）。本章的接触/速度预积分是它的\emph{同构平行实例}，读懂 \cref{ch:imu} 是读懂 \cref{sec:leg-graph} 的关键。
  \item \textbf{误差状态滤波与 InEKF}（\cref{ch:eskf}，尤其 \cref{paper:inekf}）：群仿射 + 不变误差 $\Rightarrow$ 误差传播与状态无关。本章 \cref{sec:leg-inekf} 兑现 \cref{paper:inekf} 末尾“足式接触辅助惯导”的前向许诺。
  \item \textbf{足式控制（互为表里）}（\cref{sec:ctrl-legged}）：那里用浮基动力学、接触雅可比 $\mathbf{J}_c$、摩擦锥金字塔来\emph{算力矩}；本章用\emph{同一个} $\mathbf{fk}/\mathbf{J}$ 与摩擦锥来\emph{估位姿}。估计与控制共享运动学，是足式系统“一体两面”的根源（\cref{ins:leg-control-duality}）。
\end{itemize}

\subsection*{如果跳过本章会怎样}
\begin{itemize}
  \item 你会以为足式机器人的状态估计与轮式、手持设备没区别，直接套 VIO/LIO 即可——却不知道腿提供了一路\emph{额外、近乎免费}的高频运动测量，丢掉它意味着在 IMU 单独漂移与视觉/激光退化时无所依凭；
  \item 在工程里看到 Cheetah、ANYmal、人形机器人的状态估计代码里有“contact factor”“velocity factor”“contact-aided InEKF”等模块却不知其所以然，无法解释为什么接触一切换就要重置、为什么打滑会让 $z$ 轴“鼓包”（\cref{sec:leg-open}）；
  \item 你会把 \cref{paper:inekf} 的 InEKF 当成“VIO 专用”，错过它在足式上最干净的落地——而那恰是 InEKF 当年被提出来要解决的标志性场景之一\cite{hartley2020contact}。
\end{itemize}

\begin{table}[H]\centering\small
\caption{本章阅读方式建议}
\label{tab:leg-reading}
\begin{tabular}{lll}
\toprule
阅读方式 & 时间 & 适合谁 \\
\midrule
精读（含运动学/预积分/InEKF 全推导与练习） & 4--6 小时 & 首次系统学足式状态估计、要做腿式里程计 \\
速读（跳过推导细节，看测量式与因子残差） & 1.5 小时 & 已学 IMU 预积分与 InEKF、想对齐记号 \\
速查（只看小结的符号表 + 对照表 + 速查卡） & 15 分钟 & 写代码、调参时回来查公式 \\
\bottomrule
\end{tabular}
\end{table}

\begin{note}
\textbf{本章记号与约定.}\;
沿用全书：旋转 $\mathbf{R}\in\SO(3)$、位姿 $\mathbf{T}=\big[\begin{smallmatrix}\mathbf{R}&\mathbf{t}\\\mathbf{0}^\top&1\end{smallmatrix}\big]\in\SE(3)$、\textbf{右扰动} $\mathbf{R}\,\mathrm{Exp}(\delta\boldsymbol{\phi})$、信息矩阵 $\boldsymbol{\Lambda}=\boldsymbol{\Sigma}^{-1}$。
\textbf{帧}：世界（惯性）系 $w$、机身（浮基，与 IMU 重合）系 $b$、脚（末端执行器）系 $f$、脚触地时创建的临时\textbf{接触系} $k$（刚固于地、与脚系在触地瞬间重合）。
\textbf{帧下标遵循本书约定} $\mathbf{R}_{ab}$（$b\to a$，读作“$a\leftarrow b$”），如 $\mathbf{R}_{wb}$ 把机身系映到世界系；这与 Handbook\cite{carlone2026handbook} 等用上下标 $\mathbf{R}_b^w$（$=\mathbf{R}_{wb}$）的写法互为转置约定，对照原文时须留意。本章正文统一用下标式。
\textbf{关节}：关节角 $\mathbf{q}\in\mathbb{R}^N$、关节角速度 $\dot{\mathbf{q}}$（编码器读数差分得到）。$N$ 为驱动自由度数（点足四足 $N{=}12$、平足人形 $N{=}12$ 例中各腿 $6$ DoF）。
\textbf{足端}：正运动学 $\mathbf{fk}(\mathbf{q})$、其位置/姿态分量 $\mathbf{f}_p(\mathbf{q}),\mathbf{f}_R(\mathbf{q})$，雅可比 $\mathbf{J}(\mathbf{q})$、其线/角部分 $\mathbf{J}_v,\mathbf{J}_\omega$。
\textbf{接触}：地面反力（GRF）$\mathbf{f}=[f_x,f_y,f_z]^\top$、摩擦系数 $\mu$、压力中心 CoP。
\textbf{噪声协方差记 $\boldsymbol{\Sigma}$}（避免与旋转 $\mathbf{R}$ 冲突）；编码器位置噪声 $\boldsymbol{\Sigma}_q$。本章理论主线综合 Camurri--Mattamala 的 Handbook 综述\cite{carlone2026handbook}、Bloesch 等的滤波奠基\cite{bloesch2013state}、Hartley 等的接触辅助 InEKF 与接触预积分\cite{hartley2020contact,hartley2018legged}、以及 Wisth 等的速度预积分\cite{wisth2023vilens}，并与本书 \cref{ch:imu}（预积分）、\cref{paper:inekf}（InEKF）、\cref{sec:ctrl-legged}（足式控制）逐处衔接。
\end{note}

% ════════════════════════════════════════════════════════════════
\section{为什么需要腿式里程计：从轮子到腿}
\label{sec:leg-why}

\paragraph{动机：腿是“主动悬挂”，但也是一路额外的运动测量。}
足式机器人能上楼梯、跨碎石、越障，靠的是腿把机身运动与地形剖面\emph{解耦}（一种主动悬挂）\cite{carlone2026handbook}。前几章的 SLAM 算法（视觉、激光、惯性）原样可用在足式平台上——但腿本身的关节角与关节力矩，是一路\emph{别的平台没有}的运动信息。把这路信息读成里程计，就是\textbf{腿式里程计}（leg odometry）。这个名字与轮式里程计类比而来：轮式车靠测量“轮子转了多少”推算行进距离；足式机器人则靠测量“腿相对机身怎么动了”推算机身运动。对足式\emph{控制}而言这尤其关键——闭环防摔需要高频、低漂移的实时位姿与速度，而这恰是单靠慢速外感受传感器（相机、激光）难以满足的。

\paragraph{腿与轮的根本差异：离散的接触与摆动。}
轮式车的轮子持续滚动接触地面；足式机器人却是\textbf{不断建立又断开}腿与地的接触。每条腿一个步幅分两相：\emph{摆动相}（aerial/swing，腿离地前摆）与\emph{站立相}（stance，腿与地保持\textbf{不打滑}接触）。推进机身的是站立腿。于是腿式里程计天然分解为两个子问题：
\begin{enumerate}
  \item \textbf{接触估计}（\cref{sec:leg-contact}）：判断在给定时段哪些腿处于站立相；
  \item \textbf{运动估计}（\cref{sec:leg-motion}）：由站立腿在该相内的关节信息，推算机身的增量运动。
\end{enumerate}
这条“先定站立腿、再算增量”的分解贯穿全章。注意一个用语：本章只估计机身的位姿与速度（关节状态由编码器直接测得，见 \cref{subsec:leg-sensing}），不做回环；因此本章的“状态估计”与“里程计”互换使用——里程计即\emph{没有回环的} SLAM\cite{carlone2026handbook}。

\paragraph{历史：从 Walking Truck 到商用四足与人形复兴。}
足式机器人可追溯到 $1960$ 年代试图越野的行走机械——通用电气 $1400\,\mathrm{kg}$ 的 Walking Truck 是其中代表\cite{carlone2026handbook}。$1980$ 年代末 Raibert 的单足/双足/四足跳跃机器人展示了惊人的运动能力，把腿式平台推向真实任务。最早把腿运动学用于运动估计的工作之一来自 Roston 与 Krotkov（用于 $2.5$ 吨的六足 Ambler）；此后多用粒子滤波或卡尔曼滤波把腿式里程计与 IMU、视觉融合。Bloesch 等\cite{bloesch2013state} 的里程碑工作奠定了\emph{基于滤波、全地形、运动学-惯性}的四足估计框架，并被推广到双足/人形。$2012$–$2015$ 年的 DARPA Robotics Challenge（DRC）推动了人形整机状态估计（不止 $6$-DoF 位姿，还要全身状态），各队把运动学/动力学与惯性、关节、激光\cite{carlone2026handbook}、立体视觉融合，多用卡尔曼滤波，少数探索优化与因子图。近年四足商业化又激发了更有原则、更鲁棒的方案：因子图平滑\cite{hartley2018legged,wisth2023vilens}、不变估计\cite{hartley2020contact}。$2018$–$2021$ 年的 DARPA Subterranean（SubT）挑战在极端场景下提出新难题，催生了跨腿/轮/飞行多平台的\emph{互补式}估计。$2021$ 年起，人形机器人在 AI 浪潮中复兴，把本章话题推向间歇接触（脚、躯干、手臂、手）、柔顺、长期全身估计等尚未解决的方向（\cref{sec:leg-open}）。

\begin{insight}[腿式里程计在传感器谱系中的位置：本体感受]\label{ins:leg-proprioceptive}
相机、激光是\emph{外感受}（exteroceptive）传感器——它们测的是外部环境。IMU、关节编码器、力/力矩传感器则是\emph{本体感受}（proprioceptive）——它们测的是机器人\emph{自身}的状态，不依赖外部世界是否有纹理、是否明亮、是否有结构。腿式里程计与惯性里程计同属本体感受家族：它们在黑暗、烟雾、无特征长廊里照样工作，是外感受失效时的“兜底”运动来源。这也解释了 \cref{ch:vio}/\cref{ch:imu} 反复出现的一句话——“纯惯性必败”的解药之一，正是本体感受的\emph{多路}冗余（惯性 + 腿）。腿式里程计的代价是：它要求“脚不打滑”这个关于\emph{外部世界}的隐含假设（\cref{sec:leg-contact}），一旦地面松软或打滑就失准（\cref{sec:leg-open}）。
\end{insight}

% ════════════════════════════════════════════════════════════════
\section{参考系、状态与腿运动学}
\label{sec:leg-kinematics}

本节把估计问题的“舞台”搭好：定义参考系与状态，并\textbf{自包含地}给出腿的正运动学 $\mathbf{fk}(\mathbf{q})$ 与雅可比 $\mathbf{J}(\mathbf{q})$。这两个函数是后面一切测量的来源；它们与 \cref{sec:ctrl-legged} 足式控制用的是同一套对象（\cref{ins:leg-control-duality}）。

\subsection{参考系与状态定义}
\label{subsec:leg-frames}

\paragraph{四类参考系。}
\cref{fig:leg-frames} 给出足式估计的参考系约定：世界（惯性）系 $\mathcal{F}^w$ 刚固于地，机身（浮基）系 $\mathcal{F}^b$ 刚固于机器人主躯干；不失一般性，设 IMU 与 $\mathcal{F}^b$ 重合。每个末端执行器（脚）各有一个脚系 $\mathcal{F}^f$。此外，每当一只脚触地，就\emph{临时}创建一个\textbf{接触系} $\mathcal{F}^k$：它刚固于地、垂直于地面，对人形（平足）取触地瞬间与脚系重合，对四足（点足）取与脚相同的笛卡尔位置（点足无姿态约束，可绕接触点自由转动）。接触系“静止于地”这一性质，是 \cref{sec:leg-motion} 一切推导的基石。

\begin{figure}[ht]
\centering
\begin{tikzpicture}[>=Stealth, font=\small,
  ln/.style={->, thick, gray!55!black},
  ax/.style={->, thick}]
  % world frame
  \draw[ax, third!60!black] (-3.2,-1.6) -- (-2.4,-1.6) node[right]{$x$};
  \draw[ax, third!60!black] (-3.2,-1.6) -- (-3.2,-0.8) node[above]{$z$};
  \node[third!60!black] at (-3.5,-1.9) {$\mathcal{F}^w$};
  \draw[gray!50] (-3.4,-1.6) -- (3.4,-1.6); % ground
  % body
  \draw[fill=main!10, draw=main!60!black, rounded corners] (-0.9,0.4) rectangle (0.9,1.1);
  \draw[ax, main!60!black] (0,0.75) -- (0.9,0.75) node[right]{$x$};
  \draw[ax, main!60!black] (0,0.75) -- (0,1.55) node[above]{$z$};
  \node[main!60!black] at (-0.3,1.35) {$\mathcal{F}^b$ (IMU)};
  % legs
  \draw[thick] (-0.7,0.4) -- (-1.3,-0.6) -- (-1.0,-1.6); % stance leg 1
  \draw[thick] (0.7,0.4) -- (1.2,-0.5) -- (1.5,-1.6);    % stance leg 2
  % foot/contact frame for leg1
  \fill[second!70!black] (-1.0,-1.6) circle (1.6pt);
  \draw[ax, second!70!black] (-1.0,-1.6) -- (-0.45,-1.6) node[right, font=\scriptsize]{$x$};
  \draw[ax, second!70!black] (-1.0,-1.6) -- (-1.0,-1.05) node[above, font=\scriptsize]{$z$};
  \node[second!70!black, font=\scriptsize] at (-1.35,-1.25) {$\mathcal{F}^{k}$};
  \node[font=\scriptsize] at (-1.0,-1.85) {站立 (stance)};
  % swing leg drawn lifted
  \fill[gray!50] (1.5,-1.45) circle (1.4pt);
  \node[font=\scriptsize, gray!55!black] at (1.6,-1.85) {站立/摆动};
  % T arrows
  \draw[ln, dashed] (0,0.75) to[out=-150,in=70] node[midway,left,font=\scriptsize]{$\mathbf{T}_{bf}=\mathbf{fk}(\mathbf{q})$} (-1.0,-1.45);
\end{tikzpicture}
\caption{足式机器人的参考系约定（仿 \cite{carlone2026handbook} 图 12.1 重绘）。世界系 $\mathcal{F}^w$ 固于地，机身系 $\mathcal{F}^b$ 固于躯干（与 IMU 重合）。脚触地时创建接触系 $\mathcal{F}^k$——刚固于地、垂直地面、在触地瞬间与脚系重合。正运动学 $\mathbf{fk}(\mathbf{q})$ 给出脚相对机身的位姿 $\mathbf{T}_{bf}$。}
\label{fig:leg-frames}
\end{figure}

\paragraph{状态。}
足式机器人的完整状态含机身位姿、速度与关节状态。但本章假设关节由专用传感器直接测得（\cref{subsec:leg-sensing}），故估计目标只剩机身的\textbf{位置、姿态、线速度、角速度}：
\begin{equation}
  \mathbf{x}=\begin{bmatrix}\mathbf{t}_{wb} & \mathbf{R}_{wb} & \mathbf{v} & \boldsymbol{\omega}\end{bmatrix},
  \label{eq:leg-state-raw}
\end{equation}
其中位置 $\mathbf{t}_{wb}\in\mathbb{R}^3$ 与姿态 $\mathbf{R}_{wb}\in\SO(3)$ 表机身在世界系的位姿，线/角速度 $\mathbf{v},\boldsymbol{\omega}\in\mathbb{R}^3$ 表机身的瞬时旋量。\textbf{与本书主线的接榫}：把姿态、位置、速度三件套打包，恰是 \cref{sec:se23} 的扩展位姿群 $\mathrm{SE}_2(3)$ 元素 $[\mathbf{R}_{wb},\mathbf{t}_{wb},\mathbf{v}]$；\cref{eq:leg-state-raw} 与 $\mathrm{SE}_2(3)$ 的差别仅在它额外列出角速度 $\boldsymbol{\omega}$。在装了 IMU 的实际系统里，$\boldsymbol{\omega}$ 由陀螺直接测得而不入状态，取而代之的是把 IMU 零偏增广进状态（\cref{ch:imu}）；接触辅助 InEKF 则进一步把接触点位置也增广进群（\cref{sec:leg-inekf}）。这条“先按 $\mathrm{SE}_2(3)$ 立状态、再按需增广”的脉络与 \cref{paper:inekf} 完全一致。

\subsection{腿运动学：正运动学与雅可比（自包含）}
\label{subsec:leg-fk}

\paragraph{机器人构型。}
足式机器人由一个主躯干（机身）连若干运动链（腿）构成。本章取实践中最常见两类（\cref{fig:leg-chains}）：每腿 $6$ 驱动 DoF、\emph{平足}的人形；每腿 $3$ DoF、\emph{点足}的四足。设刚体（无柔性脊柱，忽略手臂），关节为转动关节。记 $\mathbf{q}=[q_1,\dots,q_N]^\top\in\mathbb{R}^N$ 为全部 $N$ 个关节角（例中 $N{=}12$）；$\dot{\mathbf{q}}$ 为关节角速度（由编码器读数数值差分得到）。

\begin{figure}[ht]
\centering
\begin{tikzpicture}[>=Stealth, font=\scriptsize, scale=0.9]
  % quadruped
  \draw[fill=main!10, draw=main!60!black, rounded corners] (-0.7,0) rectangle (0.7,0.35);
  \node[font=\small] at (0,0.7) {四足（点足，$3$ DoF/腿）};
  \foreach \x in {-0.5,0.5} {
    \draw[thick] (\x,0) -- ($(\x,0)+(0.25,-0.5)$) -- ($(\x,0)+(0.0,-1.0)$);
    \fill[second!70!black] ($(\x,0)+(0.0,-1.0)$) circle (1.3pt);
  }
  \node at (0,-1.35) {脚 = 接触点（无姿态）};
  % humanoid
  \begin{scope}[xshift=4.6cm]
  \draw[fill=main!10, draw=main!60!black, rounded corners] (-0.45,0) rectangle (0.45,0.9);
  \node[font=\small] at (0,1.25) {人形（平足，$6$ DoF/腿）};
  \foreach \x in {-0.25,0.25} {
    \draw[thick] (\x,0) -- ($(\x,0)+(0.0,-0.5)$) -- ($(\x,0)+(0.15,-1.0)$);
    \draw[thick, line width=1.4pt] ($(\x,0)+(0.0,-1.05)$) -- ($(\x,0)+(0.35,-1.05)$);
  }
  \node at (0,-1.4) {脚 = 平面（有姿态）};
  \end{scope}
\end{tikzpicture}
\caption{两类典型构型（仿 \cite{carlone2026handbook} 图 12.2 重绘）。点足四足只能约束脚的位置（脚可绕接触点转动）；平足人形可同时约束脚的位置与姿态（脚踝提供朝向）。这一差别决定了四足用速度预积分、人形用正运动学/接触预积分（\cref{sec:leg-graph}）。}
\label{fig:leg-chains}
\end{figure}

\paragraph{正运动学 $\mathbf{fk}$：从关节角到脚位姿。}
对每只脚 $f$，定义\textbf{正运动学函数} $\mathbf{fk}(\mathbf{q}):\mathbb{R}^N\to\SE(3)$，把关节角映到脚相对机身的位姿：
\begin{equation}
  \mathbf{T}_{bf}=\mathbf{fk}(\mathbf{q})=\begin{bmatrix}\mathbf{f}_R(\mathbf{q}) & \mathbf{f}_p(\mathbf{q})\\\mathbf{0}^\top & 1\end{bmatrix}\in\SE(3),
  \label{eq:leg-fk}
\end{equation}
其中 $\mathbf{f}_R(\mathbf{q})\in\SO(3)$、$\mathbf{f}_p(\mathbf{q})\in\mathbb{R}^3$ 分别是脚在机身系中的姿态与位置。这与机械臂的正运动学是\emph{同一回事}：把腿视为一条以机身为基座的运动链，逐关节累乘相邻连杆的齐次变换即得（标准 Denavit--Hartenberg 或旋量积公式，见任一机器人学教材\cite{lynch2017modern}）；本章只需 $\mathbf{fk}$ 这个映射本身，不展开其内部连杆参数。\textbf{点足四足只用 $\mathbf{f}_p(\mathbf{q})$}：脚能绕接触点自由转动而不改变关节角，故无法从点足得到姿态测量。

\begin{insight}[估计与控制共享同一套 $\mathbf{fk}/\mathbf{J}$：足式系统的“一体两面”]\label{ins:leg-control-duality}
本章用 $\mathbf{fk}(\mathbf{q})$ 与雅可比 $\mathbf{J}(\mathbf{q})$ 把关节读数变成机身运动的\emph{测量}（估计侧）；\cref{sec:ctrl-legged} 的足式控制则用\emph{同一个} $\mathbf{J}$ 把规划出的接触力变成关节\emph{力矩}——回忆 \cref{sec:ctrl-legged-cheetah} 的 $\boldsymbol{\tau}_i=\mathbf{J}_i^\top\mathbf{R}^\top\mathbf{f}_i$（站立腿力矩 = 雅可比转置乘地面反力）。这不是巧合：雅可比 $\mathbf{J}$ 既把关节速度映到脚速度（运动学，本章），其转置 $\mathbf{J}^\top$ 又把脚处的力映到关节力矩（静力学对偶）。所以足式机器人的状态估计与运动控制\emph{共享同一个运动学内核}，互为表里：估计给控制提供机身速度与朝向（\cref{sec:ctrl-legged-cheetah} 那些量正出自本章的估计器），控制产生的接触力又反过来被本章用于接触检测（\cref{subsec:leg-grf}）。理解这层对偶，就把 P2 的感知与 P4 的控制在足式平台上接成了一个闭环。
\end{insight}

\paragraph{雅可比 $\mathbf{J}$：从关节角速度到脚速度。}
对 \cref{eq:leg-fk} 求时间导数，得\textbf{雅可比矩阵} $\mathbf{J}(\mathbf{q}):\mathbb{R}^N\to\mathbb{R}^{6\times N}$，把关节角速度映到单脚相对机身的线/角速度：
\begin{equation}
  \begin{bmatrix}\mathbf{v}_f^b\\\boldsymbol{\omega}_f^b\end{bmatrix}=\mathbf{J}(\mathbf{q})\dot{\mathbf{q}}=\begin{bmatrix}\mathbf{J}_v(\mathbf{q})\\\mathbf{J}_\omega(\mathbf{q})\end{bmatrix}\dot{\mathbf{q}},
  \label{eq:leg-jacobian}
\end{equation}
其中 $\mathbf{J}_v,\mathbf{J}_\omega\in\mathbb{R}^{3\times N}$ 是线、角部分（部分文献\cite{wensing2024legged,hartley2020contact} 把角块写在前；本章取线在前，与 \cref{eq:leg-jacobian} 一致，对照原文时留意）。由于 $\mathbf{q}$ 含所有 $N$ 个关节、而各腿运动学\emph{彼此独立}，$\mathbf{J}(\mathbf{q})$ 是\textbf{稀疏块矩阵}：只有对应该腿的列非零。记非零块为 $\bar{\mathbf{J}}(\mathbf{q})$，则第二条腿（人形）的雅可比形如
\begin{equation}
  \mathbf{J}_2(\mathbf{q})=\begin{bmatrix}\mathbf{0}_6 & \bar{\mathbf{J}}_2(\mathbf{q})\end{bmatrix},\qquad \bar{\mathbf{J}}_2(\mathbf{q})\in\mathbb{R}^{6\times6},
  \label{eq:leg-jac-block}
\end{equation}
即只有该腿那 $6$ 列非零、其余为零。这一稀疏性在 \cref{subsec:leg-grf} 从力矩反演地面反力时会被直接利用。

\begin{pitfall}[概念误区：以为腿式里程计也需要外部环境]
新手想法：“里程计不就是看着环境移动来估运动吗？腿式里程计是不是也得有相机/激光？”\textbf{错}。腿式里程计是\emph{纯本体感受}的：它只用关节角 $\mathbf{q}$、关节角速度 $\dot{\mathbf{q}}$（编码器）、IMU 与（可选的）力/接触传感器，\emph{不看}外部环境就能推算机身增量运动（\cref{ins:leg-proprioceptive}）。它对环境的唯一假设是隐含的——“站立脚不打滑”（\cref{sec:leg-contact}）。后果：若误以为腿式里程计依赖外感受，就会忽视它在黑暗/烟雾/无纹理处的独特价值，也会忽视它真正的脆弱点其实在\emph{接触假设}而非环境可见性。正确理解：腿式与惯性里程计是“兜底”的本体感受运动来源，外感受（视觉/激光）用来抑制其长期漂移并在打滑时使接触点速度可观（\cref{sec:leg-open}）。
\end{pitfall}

\subsection{关节传感：编码器、力/力矩、接触传感器}
\label{subsec:leg-sensing}

腿式里程计的输入来自三类本体感受传感器。这里只点出与里程计相关的要害，器件细节见相关文献\cite{carlone2026handbook}。
\begin{itemize}
  \item \textbf{旋转编码器}：把转轴角位置转成数字信号，是 $\mathbf{q}$ 的主要来源。分\emph{绝对式}（同一构型恒给同一读数，靠 Gray 码盘，断电不丢位）与\emph{增量式}（开机置零、靠两路 $90^\circ$ 相位的方波数增量、判转向）。绝对式因分辨率提升、成本下降正快速取代增量式。$\mathbf{q}$ 的噪声主要来自编码器量化与差分得 $\dot{\mathbf{q}}$ 的舍入误差，是 \cref{subsec:leg-encoder-noise} 不确定度的源头。
  \item \textbf{力/力矩传感器}：靠应变片测变形推算力或力矩，用于力矩控制或感知与环境的交互。$6$ 轴力/力矩传感器有时装在人形脚上直接测与地的交互。在“电机电流 $\propto$ 力矩”的反驱（backdriveable）设计下\cite{carlone2026handbook}，力矩可由电机电流间接估出，省去脚上传感器。
  \item \textbf{接触传感器}：输出“是否触地”的二值信号，常见于中小型四足的球形/圆形橡胶脚底（光学或机械式）。优点是直接给二值接触态、可硬件阈值化 $f_z$；缺点是响应慢于力/力矩传感器，且要把线缆走到脚、易受冲击损坏，故多见于小型室内四足。
\end{itemize}
关键结论：\cref{eq:leg-fk,eq:leg-jacobian} 把关节状态与机身运动联系起来，而关节状态\emph{直接}由上述传感器测得——这正是腿式里程计能把腿“读成”一路运动测量的前提。

% ════════════════════════════════════════════════════════════════
\section{运动估计：把腿读成里程计测量}
\label{sec:leg-motion}

给定关节状态与运动学，本节求机身的增量运动。有两条路：用\emph{正运动学}得两时刻间的相对位姿，或用\emph{微分运动学}瞬时估机身速度。两者都建在同一个隐含假设上——\textbf{新形成的接触系在一段时间内保持静止}。本节的两条测量式（相对位姿 \cref{eq:leg-fk-meas}、机身速度 \cref{eq:leg-vel-meas}）是后面因子（\cref{sec:leg-graph}）与滤波（\cref{sec:leg-inekf}）的共同来源。

\subsection{相对位姿估计：接触系静止的几何}
\label{subsec:leg-relpose}

\paragraph{动机：脚踩死不动时，脚相对机身的运动“反过来”就是机身相对地的运动。}
设机器人在 $zx$ 平面内行走（\cref{fig:leg-relpose}），机身在相邻两时刻取位姿 $\mathcal{F}^b$ 与 $\mathcal{F}^{b'}$，对应同一只站立脚的脚系与关节角分别为 $(\mathcal{F}^f,\mathbf{q})$ 与 $(\mathcal{F}^{f'},\mathbf{q}')$。因接触系 $\mathcal{F}^k$ 静止于地、且在站立期内脚系与接触系重合，机身向前移动多少，脚（在机身系看）就向后移动多少。形式化：

\begin{figure}[ht]
\centering
\begin{tikzpicture}[>=Stealth, font=\scriptsize,
  ax/.style={->, semithick}]
  \draw[gray!50] (-0.5,-1.7) -- (5.3,-1.7); % ground
  % contact point (shared, stationary)
  \fill[second!70!black] (2.4,-1.7) circle (1.6pt);
  \node[second!70!black] at (2.4,-2.0) {$\mathcal{F}^{k}$ (静止)};
  % body at time 1
  \draw[fill=main!10, draw=main!60!black, rounded corners] (0.4,0.4) rectangle (1.6,0.85);
  \draw[ax, main!60!black] (1.0,0.62) -- (1.6,0.62); \draw[ax, main!60!black] (1.0,0.62) -- (1.0,1.2);
  \node[main!60!black] at (0.7,1.05) {$\mathcal{F}^{b}$};
  \draw[thick] (1.0,0.4) -- (1.7,-0.6) -- (2.4,-1.7);
  % body at time 2 (moved forward)
  \draw[fill=main!6, draw=main!50!black, rounded corners] (2.6,0.4) rectangle (3.8,0.85);
  \draw[ax, main!50!black] (3.2,0.62) -- (3.8,0.62); \draw[ax, main!50!black] (3.2,0.62) -- (3.2,1.2);
  \node[main!50!black] at (3.6,1.05) {$\mathcal{F}^{b'}$};
  \draw[thick, gray!60] (3.2,0.4) -- (2.9,-0.6) -- (2.4,-1.7);
  % displacement arrow
  \draw[ax, third!60!black, dashed] (1.0,0.62) -- node[above]{$\mathbf{t}_{b b'}$（机身前进）} (3.2,0.62);
  \draw[ax, second!70!black, dashed] (2.4,-1.55) to[out=120,in=-120] node[left,xshift=-1mm]{脚后退} (1.7,-0.6);
\end{tikzpicture}
\caption{接触系静止时的相对位姿（仿 \cite{carlone2026handbook} 图 12.5 重绘）。机身从 $\mathcal{F}^b$ 前进到 $\mathcal{F}^{b'}$，与此同时同一只站立脚（在机身系看）向后移动等量。由 $\mathbf{fk}$ 在两时刻的值即可读出机身相对位姿，无需外部环境。}
\label{fig:leg-relpose}
\end{figure}

\begin{derivation}[由正运动学得相对位姿]
接触系静止意味着 $\mathbf{T}_{kb}$ 与 $\mathbf{T}_{kb'}$ 共享同一固定的 $\mathcal{F}^k$。机身从 $b$ 到 $b'$ 的相对位姿 $\mathbf{T}_{bb'}$（把 $b'$ 系坐标变到 $b$ 系）可经接触系中转：
\[
  \mathbf{T}_{bb'}=\mathbf{T}_{bk}\,\mathbf{T}_{kb'}=\mathbf{T}_{bk}(\mathbf{T}_{b'k})^{-1}.
\]
站立期内脚系与接触系重合（$\mathcal{F}^k\equiv\mathcal{F}^f$），故 $\mathbf{T}_{bk}=\mathbf{T}_{bf}=\mathbf{fk}(\mathbf{q})$、$\mathbf{T}_{b'k}=\mathbf{T}_{b'f'}=\mathbf{fk}(\mathbf{q}')$。代入得机身\textbf{相对位姿测量}：
\begin{equation}
  \mathbf{T}_{bb'}=\mathbf{fk}(\mathbf{q})\,\mathbf{fk}(\mathbf{q}')^{-1}
  \quad\Longleftrightarrow\quad
  \mathbf{T}_{b'b}=\mathbf{fk}(\mathbf{q}')^{-1}\,\mathbf{fk}(\mathbf{q}),
  \label{eq:leg-fk-meas}
\end{equation}
两式互为逆，分别表“$b'$ 系在 $b$ 系”与“$b$ 系在 $b'$ 系”的位姿。
\end{derivation}

\cref{eq:leg-fk-meas} 在关节状态与机身相对位姿之间建了一座桥：\textbf{把相对位姿一段段串起来，就得到纯靠关节读数的航位推算}（dead reckoning）。但它\emph{只}在站立期成立；且对点足机器人，要靠它定出 $6$ 自由度相对位姿需\emph{至少三只脚同时站立}——这对很多步态（如奔驰中四脚同时离地的腾空相）并不现实。

\paragraph{增广接触点：四足与人形的状态。}
为绕开“至少三足站立”的限制，四足的标准做法\cite{bloesch2013state} 是把状态 \cref{eq:leg-state-raw} \textbf{增广}上各腿接触系在世界系的位置 $\mathbf{c}_i\in\mathbb{R}^3$；人形（有脚踝、平足）还可加接触点姿态 $\mathbf{B}_i\in\SO(3)$。又因任意时刻未必有足够腿在站立（腾空相），常假设装有 IMU——于是状态里不再含角速度 $\boldsymbol{\omega}$（陀螺直接测），转而纳入 IMU 零偏 $\mathbf{b}^a,\mathbf{b}^g$（\cref{ch:imu}）。综合得四足与人形的状态：
\begin{align}
  \mathbf{x}^{\text{quad}}&=\begin{bmatrix}\mathbf{t} & \mathbf{R} & \mathbf{v} & \mathbf{b}^a & \mathbf{b}^g & \mathbf{c}_1 & \mathbf{c}_2 & \mathbf{c}_3 & \mathbf{c}_4\end{bmatrix},\label{eq:leg-state-quad}\\
  \mathbf{x}^{\text{hum}}&=\begin{bmatrix}\mathbf{t} & \mathbf{R} & \mathbf{v} & \mathbf{b}^a & \mathbf{b}^g & \mathbf{c}_1 & \mathbf{c}_2 & \mathbf{B}_1 & \mathbf{B}_2\end{bmatrix}.\label{eq:leg-state-hum}
\end{align}
有了显式的接触位姿，\cref{eq:leg-fk-meas} 对任意第 $i$ 腿可写成一条把机身位姿、接触位姿与正运动学联系起来的约束。记机身位姿 $\mathbf{T}=\big[\begin{smallmatrix}\mathbf{R}&\mathbf{t}\\\mathbf{0}^\top&1\end{smallmatrix}\big]$（世界系）、第 $i$ 腿接触位姿 $\mathbf{C}_i=\big[\begin{smallmatrix}\mathbf{B}_i&\mathbf{c}_i\\\mathbf{0}^\top&1\end{smallmatrix}\big]$（世界系），则脚相对机身的位姿 $\mathbf{T}_{bf}=\mathbf{fk}(\mathbf{q})=\mathbf{T}^{-1}\mathbf{C}_i$。展开：
\begin{equation}
  \mathbf{T}^{-1}\mathbf{C}_i=\begin{bmatrix}\mathbf{R}^\top & -\mathbf{R}^\top\mathbf{t}\\\mathbf{0}^\top & 1\end{bmatrix}\begin{bmatrix}\mathbf{B}_i & \mathbf{c}_i\\\mathbf{0}^\top & 1\end{bmatrix}=\begin{bmatrix}\mathbf{R}^\top\mathbf{B}_i & \mathbf{R}^\top(\mathbf{c}_i-\mathbf{t})\\\mathbf{0}^\top & 1\end{bmatrix},
  \label{eq:leg-Tc-expand}
\end{equation}
对照 \cref{eq:leg-fk} 的上三角块，读出正运动学的位置/姿态部分：
\begin{equation}
  \mathbf{f}_p(\mathbf{q})=\mathbf{R}^\top(\mathbf{c}_i-\mathbf{t}),\qquad \mathbf{f}_R(\mathbf{q})=\mathbf{R}^\top\mathbf{B}_i.
  \label{eq:leg-fk-pos-rot}
\end{equation}
\cref{eq:leg-fk-pos-rot} 即用作估计框架（滤波或因子图）里的\textbf{腿式里程计观测}——它把\emph{可测}的左端（正运动学，来自编码器）与\emph{待估}的右端（机身位姿 + 接触位姿）联系起来。点足四足只用位置式（无法从点足测姿态）。

\begin{pitfall}[概念误区：把 $\mathbf{c}_i$ 当成一个永久固定的“地标”]
新手想法：“脚一旦落地，$\mathbf{c}_i$ 就是个不动的世界点，跟视觉路标一样永久。”\textbf{不全对}。$\mathbf{c}_i$ 只在该脚\emph{当前这次站立期}内固定；脚一抬起再落下，是一次\emph{新}的接触，$\mathbf{c}_i$ 必须\textbf{重置}为新落点的位置（并放大其不确定度，因为新落点的世界坐标是靠刚刚估出的机身位姿外推的）。后果：若不在接触切换时重置 $\mathbf{c}_i$，滤波器会把“脚抬起后的摆动轨迹”错当成机身运动，瞬间发散。正确做法：接触检测（\cref{sec:leg-contact}）一旦判定某脚由摆动转站立，就用当前机身位姿 + 正运动学初始化该腿的 $\mathbf{c}_i$，并在接触保持期内才信任与它相关的测量。接触预积分因子（\cref{subsec:leg-contact-preint}）的“只在同一站立腿、同一站立相内有效”正源于此。
\end{pitfall}

\subsection{速度估计：接触点零速度约束}
\label{subsec:leg-velocity}

\paragraph{动机：不想维护接触点状态，就直接估机身速度。}
\cref{eq:leg-fk-meas} 要么需三足站立、要么需把接触位姿纳入状态。微分运动学给出另一条路：从每只站立腿\emph{瞬时}估出机身\textbf{速度}。这在四足上广为采用\cite{bloesch2013state}，人形上较少。优点是速度测量易（预）积分进滤波或因子、不留历史故不累积位置误差、无需维护额外状态；代价是 $\dot{\mathbf{q}}$ 多由 $\mathbf{q}$ 数值差分得到，舍入误差可能拖累精度。

\begin{derivation}[由零速度约束得机身速度]
核心观察：站立腿的接触点 $k$ 必须\emph{静止于世界系}（不打滑），故其在世界系的速度为零：
\begin{equation}
  \mathbf{v}_k^w=\mathbf{0}.
  \label{eq:leg-zero-vel}
\end{equation}
同时，接触点速度（在机身系）等于脚速度，由雅可比 \cref{eq:leg-jacobian} 给出：
\begin{equation}
  \mathbf{v}_k^b=\mathbf{v}_f^b=\mathbf{J}_v(\mathbf{q})\dot{\mathbf{q}}.
  \label{eq:leg-vk-body}
\end{equation}
把接触点世界速度按刚体运动学展开（机身速度 + 角速度叉乘杆臂 + 接触点相对机身速度），并代入零速度约束 \cref{eq:leg-zero-vel}：
\[
  \underbrace{\mathbf{v}_k^w}_{=\,\mathbf{0}}=\mathbf{v}_b^w+\boldsymbol{\omega}_b^w\times(\mathbf{R}_{wb}\mathbf{t}_k^b)+\mathbf{R}_{wb}\,\mathbf{v}_k^b,
\]
其中 $\mathbf{t}_k^b=\mathbf{f}_p(\mathbf{q})$ 是机身到接触点的杆臂、$\mathbf{v}_k^b=\mathbf{J}_v(\mathbf{q})\dot{\mathbf{q}}$。整理出\textbf{机身（世界系）速度}：
\begin{equation}
  \boxed{\;\mathbf{v}_b^w=-\boldsymbol{\omega}_b^w\times\big(\mathbf{R}_{wb}\,\mathbf{f}_p(\mathbf{q})\big)-\mathbf{R}_{wb}\,\mathbf{J}_v(\mathbf{q})\dot{\mathbf{q}}\;}
  \label{eq:leg-vel-meas}
\end{equation}
（同一速度若改在\emph{机身系}表达，对整式左乘 $\mathbf{R}_{wb}^\top$：$\mathbf{v}_b^b=\mathbf{R}_{wb}^\top\mathbf{v}_b^w=-\boldsymbol{\omega}_b^b\times\mathbf{f}_p(\mathbf{q})-\mathbf{J}_v(\mathbf{q})\dot{\mathbf{q}}$，与 Handbook 式（12.21）一致。注意机身系式里杆臂 $\mathbf{f}_p$ \emph{不再}需要旋转——因角速度 $\boldsymbol{\omega}_b^b$ 与杆臂 $\mathbf{f}_p$ 同在机身系；而 \cref{eq:leg-vel-meas} 世界系式里 $\boldsymbol{\omega}_b^w$ 与机身系杆臂不同系，须先用 $\mathbf{R}_{wb}$ 把杆臂转到世界系再叉乘。两式是同一物理速度在两坐标系的表达，差一个\emph{整体}旋转 $\mathbf{R}_{wb}$，并非逐项增删。）
\end{derivation}

\cref{eq:leg-vel-meas} 把机身绝对速度写成正运动学 $\mathbf{f}_p$ 与微分运动学 $\mathbf{J}_v$ 的函数，可直接作滤波的测量更新或因子图的速度因子（\cref{subsec:leg-vel-preint}）。其中角速度 $\boldsymbol{\omega}_b^w$ 由 IMU 陀螺测得（故只需解线速度），杆臂项 $-\boldsymbol{\omega}_b^w\times(\mathbf{R}_{wb}\mathbf{f}_p)$ 正是 \cref{ins:imu-leverarm} 中“传感器不在体心”那个杆臂效应在足端的体现——脚离机身越远、机身转得越快，这一项越不可忽略。

\begin{insight}[两条测量、两类平台：相对位姿 vs 速度]\label{ins:leg-two-measurements}
\cref{eq:leg-fk-meas}（相对位姿）与 \cref{eq:leg-vel-meas}（速度）是腿式里程计的两块基石，分工与平台相关：
\begin{itemize}
  \item \textbf{人形（平足）}：脚踝给出脚的\emph{朝向}，故能用正运动学约束完整 $6$ 自由度相对位姿 $\Rightarrow$ 走\emph{正运动学因子 + 接触预积分}（\cref{subsec:leg-fk-factor,subsec:leg-contact-preint}）。
  \item \textbf{四足（点足）}：脚能绕接触点自由转动，正运动学\emph{无法}约束 $6$ 自由度相对位姿；但零速度约束仍给出机身\emph{线}速度 $\Rightarrow$ 走\emph{速度预积分}（\cref{subsec:leg-vel-preint}）。
\end{itemize}
这就是为什么后面因子图要分两套写法——不是数学上有本质不同，而是\emph{点足缺一维姿态信息}。两者的统一视角在 \cref{sec:leg-inekf}：无论哪种，接触约束都可写进 $\mathrm{SE}_2(3)$（或其增广）上的不变滤波。
\end{insight}

至此都假设站立腿已知。下一节补上“如何判定站立腿”。

% ════════════════════════════════════════════════════════════════
\section{接触估计：哪条腿在站立}
\label{sec:leg-contact}

接触估计的定义随应用而变。对协作机械臂，末端“碰到东西”（有不可忽略外力）即算接触；但对腿式里程计，一只脚\emph{只有当接触点随时间静止}（点足即不打滑）时才算站立。本节先给“不打滑”的力学条件，再说明为何工程上普遍\textbf{只阈值化法向力 $f_z$}，最后讲无专用传感器时如何从关节力矩反演地面反力。

\subsection{不打滑条件：摩擦圆锥与压力中心}
\label{subsec:leg-friction}

\paragraph{动机：站立 = 接触点不动 = 力落在摩擦锥内。}
脚不打滑，等价于地面反力（GRF）$\mathbf{f}=[f_x,f_y,f_z]^\top$ 落在\textbf{摩擦圆锥}内——切向力不超过法向力乘摩擦系数：
\begin{equation}
  \sqrt{f_x^2+f_y^2}\le\mu_{xy}\,f_z,
  \label{eq:leg-friction-cone}
\end{equation}
其中 $f_x,f_y$ 是相对接触平面的切向分量（依地形局部形态而定），$\mu_{xy}$ 是摩擦系数（依地面与脚的材质而定）。\textbf{与足式控制的呼应}：这正是 \cref{sec:ctrl-legged-contact} 控制侧用的同一个摩擦圆锥 $\|\boldsymbol{\lambda}_t\|\le\mu\lambda_n$——控制端把它\emph{线性化成方形金字塔}塞进 QP 以\emph{保证}规划出的接触力不打滑（\cref{eq:ctrl-friction-pyramid} 那段），估计端则用它\emph{判定}当前脚是否在站立。一个约束，控制用来“不让打滑”，估计用来“判断有没有打滑”，又一处“一体两面”（\cref{ins:leg-control-duality}）。

\paragraph{平足还要不转动：压力中心。}
人形平足还能对地施加\emph{力矩}，于是除 \cref{eq:leg-friction-cone} 外还需脚不绕竖轴转、不翻边。前者要切向摩擦力矩受限、后者要压力中心（Center of Pressure, CoP）落在支撑多边形内：
\begin{equation}
  \begin{bmatrix}-\tau_y/f_z\\\tau_x/f_z\end{bmatrix}\le\begin{bmatrix}\mathrm{CoP}_x\\\mathrm{CoP}_y\end{bmatrix},\qquad |\tau_z|\le\mu_z f_z,
  \label{eq:leg-cop}
\end{equation}
其中 $\boldsymbol{\tau}$ 是接触力矩，$\mu_z$ 是转动摩擦系数，$\mathrm{CoP}_{x},\mathrm{CoP}_{y}$ 是压力中心分量的上界（由接触面几何决定的支撑多边形边界）。\cref{eq:leg-cop} 的几何含义与 \cref{sec:ctrl-legged} 控制中“接触力矩不出支撑域”一致——支撑多边形越大（脚越大），平足人形越不易翻倒。

\subsection{为何只阈值化 $f_z$}
\label{subsec:leg-fz}

\paragraph{动机：一个足够大的法向力，自动满足所有不打滑条件。}
\cref{eq:leg-friction-cone,eq:leg-cop} 的右端都含 $f_z$。法向力 $f_z$ 足够大时，无论切向力与力矩取何值，这些不等式都更容易满足——直觉上“踩得够实”就不会滑、不会翻。因此实践中最常用的接触估计就是\textbf{对 $f_z$ 设阈值}：$f_z$ 超过某门限即判该脚站立。不同实现的差别仅在\emph{$f_z$ 怎么测/估}与机器人的具体特性：
\begin{itemize}
  \item \textbf{接触传感器}（\cref{subsec:leg-sensing}）：硬件层面隐式阈值化 $f_z$——二值信号只在力超额定值时触发，腿式里程计直接信这个二值态，最简单。
  \item \textbf{力/力矩传感器}：直接测得随时间变化的 $f_z$，可把力的\emph{模式}对应到非二值事件。例如：小而上升、持续一段的力 = 脚正在着地但尚未站稳；力降到某值（如单腿额定载荷一半）以下 = 脚即将离地、不再可信。两种情形都应丢弃该腿测量或放大其不确定度\cite{carlone2026handbook}。
  \item \textbf{IMU（脚/腿上）}：脚着地的冲击会改变脚的加速度，故腿上 IMU 可\emph{隐式}检测站立腿。好处是传感器不直接承受冲击、不易坏，代价是需额外信号处理来识别加速度突变。
  \item \textbf{关节力矩}：无专用传感器时，从关节力矩 + 机器人动力学反演 GRF（下节）。
\end{itemize}

\begin{pitfall}[工程陷阱：固定 $f_z$ 阈值在不同步态/地形上失效]
错误做法：标定时在平地慢走测一个 $f_z$ 阈值，全场景照用。现象：快步/跳跃时着地冲击使 $f_z$ 瞬间冲高、误判摆动腿为站立；软地/下坡时正常站立的 $f_z$ 偏低、漏判站立腿。根本原因：$f_z$ 的量级随步态动力学（着地速度）、地形坡度、负载显著变化，单一固定门限刻画不了。正确做法：用滞回（双门限，进/出站立用不同阈值防抖）、结合力的\emph{时间模式}（上升沿/持续时长，而非瞬时值）、或概率化地融合运动学与动力学线索\cite{carlone2026handbook}；更现代的做法是用学习的接触分类器（\cref{sec:leg-open}）。自检：把 $f_z$ 曲线与真值接触态（如用动捕或外部力板标注）对齐，看固定阈值的误判率随步态变化多大。
\end{pitfall}

\subsection{从关节力矩反演地面反力}
\label{subsec:leg-grf}

\paragraph{动机：脚上加传感器并非总可行。}
不同机器人的形态、设计、集成约束未必允许在脚上加力/接触传感器。此时可借浮基动力学，从关节力矩 $\boldsymbol{\tau}$ 估出 GRF。浮基系统的动力学（用递推 Newton--Euler 算\cite{featherstone2008rigid}，与 \cref{sec:ctrl-legged} 控制侧用的是同一套浮基动力学方程）把质量阵、Coriolis/离心/重力项、接触雅可比与力矩联系起来；利用 $\mathbf{J}_q$（站立腿雅可比的堆叠）的块状结构，对四足可逐腿解出脚处 GRF：
\begin{equation}
  \mathbf{f}_i=-(\bar{\mathbf{J}}_{i,v}^\top)^{-1}\big(\boldsymbol{\tau}_i-\mathbf{h}_{q,i}-\mathbf{F}^\top\dot{\mathbf{v}}\big),
  \label{eq:leg-grf-from-torque}
\end{equation}
其中 $\mathbf{f}_i,\boldsymbol{\tau}_i\in\mathbb{R}^3$ 是第 $i$ 腿的 GRF 与关节力矩，$\bar{\mathbf{J}}_{i,v}$ 是该脚线雅可比的非零块（四足为方阵、可逆），$\mathbf{F}$ 是质量阵的一块，$\mathbf{h}_{q,i}$ 是该腿的 Coriolis/离心/重力项。

\paragraph{两个缺失的量。}
\cref{eq:leg-grf-from-torque} 给出的 $\mathbf{f}_i$ 只能在机身系——要还原\emph{真实} GRF 还差两条信息：
\begin{enumerate}
  \item \textbf{地形局部坡度}（接触力的朝向依赖它）：人形脚踝能较好近似；四足无外感受时定不出接触系朝向，只能从其他站立脚启发式推断（如对它们拟合一个平面）。
  \item \textbf{摩擦系数}（切向力依赖它）：依地面材质而定，只能先验给定或从打滑量反推\cite{carlone2026handbook}。
\end{enumerate}
总的来说，单凭关节传感判定接触态仍是\emph{开放问题}：已有概率法（融合运动学与动力学）与学习法（\cref{sec:leg-open}）。

\begin{insight}[接触估计是腿式里程计的“数据关联”]\label{ins:leg-contact-assoc}
回想 \cref{ch:vo}/\cref{ch:lidar_slam}：视觉/激光里程计的精度上限往往不在几何求解，而在\emph{数据关联}是否正确（匹配错一个特征，整帧就偏）。腿式里程计的“数据关联”就是接触估计：哪条腿在站立、从何时到何时——判错一次，\cref{eq:leg-fk-meas}/\cref{eq:leg-vel-meas} 就把摆动腿的运动错当机身运动，注入一个大粗差。所以接触估计虽看似“辅助”，却与运动学求解\emph{同等重要}；这也是为什么本章把它单列一节，并在 \cref{sec:leg-open} 把“学习接触估计”列为头号前沿——正如视觉里程计把鲁棒匹配（\cref{ch:vo} 的 RANSAC、鲁棒核）当作命门。
\end{insight}

% ════════════════════════════════════════════════════════════════
\section{因子图平滑：腿式预积分}
\label{sec:leg-graph}

有了腿式测量（相对位姿 \cref{eq:leg-fk-meas}、机身速度 \cref{eq:leg-vel-meas}）与站立腿判定（\cref{sec:leg-contact}），本节把它们送进估计框架。历史上足式估计为满足闭环控制的高频需求，多用卡尔曼滤波的非线性变体（EKF\cite{bloesch2013state}、UKF、InEKF\cite{hartley2020contact}），融合惯性与运动学喂给控制器；因子图平滑近几年才被采用，且多用于较慢的外感受传感器（激光/相机）做建图\cite{hartley2018legged,wisth2023vilens}。无论哪种，要把腿式测量\emph{最优}地融进去，都得先量化其不确定度——这是本节起点；然后给出腿式因子的两套预积分写法。这两套因子与雷达、事件因子一样，都不外乎 \cref{ins:radar-factor-recipe} 的“三问入图”（约束哪些状态、残差怎么写、信息矩阵 $\boldsymbol\Lambda$ 多大），只是把测量从 IMU 换成了编码器与接触运动学。

\subsection{编码器噪声的一阶传播}
\label{subsec:leg-encoder-noise}

\paragraph{动机：腿式测量的不确定度根在编码器。}
腿式里程计的主要噪声源是关节编码器。把它建成加性零均值高斯：真值 $\mathbf{q}$ 与测量 $\tilde{\mathbf{q}}$ 相差 $\boldsymbol{\eta}_q\sim\mathcal{N}(\mathbf{0},\boldsymbol{\Sigma}_q)$：
\begin{equation}
  \tilde{\mathbf{q}}=\mathbf{q}+\boldsymbol{\eta}_q.
  \label{eq:leg-encoder-noise}
\end{equation}
正运动学与微分运动学含旋转，是\emph{非线性}的，不保高斯。但与 \cref{ch:imu} 完全同套的招式——用雅可比做\textbf{一阶近似}（局部线性、保高斯）——即可把编码器噪声传到测量：
\begin{equation}
  \mathbf{f}_p(\mathbf{q}+\boldsymbol{\eta}_q)\approx\mathbf{f}_p(\mathbf{q})+\mathbf{J}_c(\mathbf{q})\,\boldsymbol{\eta}_q,
  \label{eq:leg-fp-noise}
\end{equation}
其中 $\mathbf{J}_c(\mathbf{q})$ 是\emph{接触系下}表达的机身雅可比（即 $\mathbf{J}(\mathbf{q})$ 但表在接触系）。对速度测量，编码器位置噪声 $\boldsymbol{\eta}_q$ 与角速度噪声 $\boldsymbol{\eta}_{\dot{q}}$ 同时作用\cite{wisth2023vilens}：
\begin{equation}
  \mathbf{J}(\mathbf{q}+\boldsymbol{\eta}_q)(\dot{\mathbf{q}}+\boldsymbol{\eta}_{\dot{q}})\approx\mathbf{J}(\mathbf{q})\dot{\mathbf{q}}+\frac{\partial}{\partial\mathbf{q}}\big(\mathbf{J}(\mathbf{q})\dot{\mathbf{q}}\big)\boldsymbol{\eta}_q+\mathbf{J}(\mathbf{q})\boldsymbol{\eta}_{\dot{q}}.
  \label{eq:leg-vel-noise}
\end{equation}
\cref{eq:leg-fp-noise,eq:leg-vel-noise} 中乘在噪声上的项，都可并成一个对给定编码器测量\emph{已知}的线性系数——于是腿式测量的协方差就由 $\boldsymbol{\Sigma}_q$（与 $\boldsymbol{\Sigma}_{\dot{q}}$）经这些雅可比线性传播而来。这与 \cref{ch:imu} 把 IMU 原始噪声经 $\mathbf{A},\mathbf{B}$ 矩阵传到预积分协方差（\cref{eq:cov-recursion}）是\emph{同一种线性化传播思想}，只是这里的“原始噪声”是编码器、传播算子是运动学雅可比。

\subsection{两个例子的总览}
\label{subsec:leg-two-examples}

下面给两个文献里的代表例，展示\textbf{预积分理论}（\cref{ch:imu}）与本章腿式测量如何在因子图里结合：人形的\emph{接触预积分}\cite{hartley2018legged} 与四足的\emph{速度预积分}\cite{wisth2023vilens}。两例都把 \cref{sec:leg-motion} 的测量重写成因子的\textbf{残差与协方差}。先把因子图的共同骨架立起来（\cref{fig:leg-factor-graph}）：一个先验因子锚定图、一个\textbf{预积分 IMU 因子}（\cref{ch:imu}）引入惯性运动先验，再叠加腿式因子；外感受测量（激光/相机）可作额外因子随手加入。

\begin{figure}[ht]
\centering
\begin{tikzpicture}[>=Stealth, font=\scriptsize,
  state/.style={circle, draw, minimum size=7mm, inner sep=0pt, fill=main!10, draw=main!60!black},
  fac/.style={rectangle, draw, minimum size=2.4mm, inner sep=0pt},
  imuf/.style={fac, fill=orange!75!red},
  fkf/.style={fac, fill=second!70!black},
  conf/.style={fac, fill=blue!55!black},
  velf/.style={fac, fill=teal!65!black},
  extf/.style={fac, fill=magenta!70!black},
  prf/.style={fac, fill=black},
  ln/.style={thick, gray!55!black}]
  \node[state] (xi) {$\mathbf{x}_i$};
  \node[state, right=22mm of xi] (xj) {$\mathbf{x}_j$};
  \node[state, right=22mm of xj] (xk) {$\mathbf{x}_k$};
  % prior
  \node[prf, left=5mm of xi] (pr) {}; \draw[ln] (pr)--(xi);
  % imu factors
  \node[imuf] (imu1) at ($(xi)!0.5!(xj)$) {}; \draw[ln] (xi)--(imu1)--(xj);
  \node[imuf] (imu2) at ($(xj)!0.5!(xk)$) {}; \draw[ln] (xj)--(imu2)--(xk);
  % leg factors (below)
  \node[velf] (v1) at ($(xi)!0.5!(xj)+(0,-7mm)$) {}; \draw[ln] (xi) to[bend right=18] (v1); \draw[ln] (v1) to[bend right=18] (xj);
  \node[velf] (v2) at ($(xj)!0.5!(xk)+(0,-7mm)$) {}; \draw[ln] (xj) to[bend right=18] (v2); \draw[ln] (v2) to[bend right=18] (xk);
  % exteroceptive
  \node[extf] (e1) at ($(xi)!0.5!(xk)+(0,9mm)$) {}; \draw[ln, magenta!60!black] (xi) to[bend left=22] (e1); \draw[ln, magenta!60!black] (e1) to[bend left=22] (xk);
  % legend
  \node[font=\scriptsize, anchor=west] at ($(xk)+(8mm,4mm)$) {\tikz\node[prf]{};\, 先验};
  \node[font=\scriptsize, anchor=west] at ($(xk)+(8mm,1.5mm)$) {\tikz\node[imuf]{};\, IMU 预积分};
  \node[font=\scriptsize, anchor=west] at ($(xk)+(8mm,-1mm)$) {\tikz\node[velf]{};\, 腿式（接触/速度）};
  \node[font=\scriptsize, anchor=west] at ($(xk)+(8mm,-3.5mm)$) {\tikz\node[extf]{};\, 外感受（激光/相机）};
\end{tikzpicture}
\caption{腿式估计的因子图骨架（仿 \cite{carlone2026handbook} 图 12.7、\cite{hartley2018legged,wisth2023vilens} 重绘）。关键帧状态 $\mathbf{x}_i$（圆）由先验、预积分 IMU 因子（\cref{ch:imu}）、腿式因子（接触或速度预积分）约束；外感受测量（约束两状态）可作额外因子加入——与 \cref{ch:vio} 的 VIO、\cref{ch:lidar_slam} 的 LIO-SAM 因子图是同一种结构，只是多了一类腿式因子。}
\label{fig:leg-factor-graph}
\end{figure}

\subsection{正运动学因子（人形）}
\label{subsec:leg-fk-factor}

正运动学因子把脚接触系的位姿、机身位姿、与站立腿正运动学三者绑在一起（都在世界系）。把编码器噪声 \cref{eq:leg-fp-noise} 代入相对位姿测量 \cref{eq:leg-fk-pos-rot}，定义\textbf{正运动学因子}的残差与协方差：
\begin{align}
  \mathbf{r}_{\mathcal{F}}&=\mathrm{Log}\!\big(\mathbf{C}_i^{-1}\,\mathbf{T}\,\mathbf{fk}(\tilde{\mathbf{q}})\big),\label{eq:leg-fk-residual}\\
  \boldsymbol{\Sigma}_{\mathcal{F}}&=\mathbf{J}_c(\tilde{\mathbf{q}})\,\boldsymbol{\Sigma}_q\,\mathbf{J}_c^\top(\tilde{\mathbf{q}}),\label{eq:leg-fk-cov}
\end{align}
其中 $\mathbf{T}$ 是机身位姿、$\mathbf{C}_i$ 是第 $i$ 腿接触位姿（\cref{eq:leg-state-hum}）。残差 $\mathbf{r}_{\mathcal{F}}$ 度量“接触位姿与机身位姿之差”偏离正运动学预测的程度，权重为编码器噪声传播得到的 $\boldsymbol{\Sigma}_{\mathcal{F}}^{-1}$。这与 \cref{ch:imu} 的预积分 IMU 因子残差 \cref{eq:imu-residual} 同型：都是“测量预测 $\ominus$ 状态给出的量”经 $\mathrm{Log}$ 映到切空间、再用协方差逆加权。

\subsection{接触预积分因子（人形）}
\label{subsec:leg-contact-preint}

\paragraph{动机：把站立期内接触点的“缓慢漂移”预积分成一条因子。}
接触预积分因子\cite{hartley2018legged} 在接触点上加约束：理想（不打滑）时接触位姿应\emph{不变}；实际有打滑，建成加在接触点速度上的高斯噪声。该因子刻画两时刻间接触点因这一噪声能漂多少——正是 \cref{ch:imu} 预积分思想的搬用：\textbf{把站立期内的高频接触状态积成与两端状态无关的相对增量，再作一条因子}。

\begin{derivation}[接触预积分残差]
设两相邻状态 $\mathbf{x}_i,\mathbf{x}_j$（时刻 $t_i,t_j$），用接触位姿 $\mathbf{C}_i,\mathbf{C}_j$ 的旋转/平移分量 $(\mathbf{B},\mathbf{c})$。引入预积分接触噪声 $\delta\boldsymbol{\theta}_{ij}$（角）、$\delta\mathbf{d}_{ij}$（平移），把“理想不变”写成
\begin{align}
  \Delta\tilde{\mathbf{B}}_{ij}&=\mathbf{B}_i^\top\mathbf{B}_j\,\mathrm{Exp}(\delta\boldsymbol{\theta}_{ij})=\mathbf{I},\label{eq:leg-contact-rot}\\
  \Delta\tilde{\mathbf{c}}_{ij}&=\mathbf{B}_i^\top(\mathbf{c}_j-\mathbf{c}_i)+\delta\mathbf{d}_{ij}=\mathbf{0},\label{eq:leg-contact-trans}
\end{align}
即“接触系姿态不变、位置不变”各被一个零均值高斯噪声扰动。把 \cref{eq:leg-contact-rot,eq:leg-contact-trans} 重排堆叠成一个向量残差，并把噪声移到末尾，得\textbf{接触预积分因子}：
\begin{equation}
  \mathbf{r}_{\mathcal{C}}=\begin{bmatrix}\mathrm{Log}\!\big(\mathbf{B}_i^\top\mathbf{B}_j\big)\\\mathbf{B}_i^\top(\mathbf{c}_j-\mathbf{c}_i)\end{bmatrix},\qquad
  \boldsymbol{\Sigma}_{\mathcal{C}}=\begin{bmatrix}\boldsymbol{\Sigma}_w & \mathbf{0}\\\mathbf{0} & \boldsymbol{\Sigma}_v\end{bmatrix}\Delta t_{ij},
  \label{eq:leg-contact-preint}
\end{equation}
其中协方差 $\boldsymbol{\Sigma}_{\mathcal{C}}$ 由接触角速度协方差 $\boldsymbol{\Sigma}_w$ 与接触线速度协方差 $\boldsymbol{\Sigma}_v$ 在 $\Delta t_{ij}$ 上\emph{时间积分}得到（站立越久、积累的允许漂移越大）。
\end{derivation}

\begin{insight}[接触预积分 = IMU 预积分换一种“高频测量”]\label{ins:leg-preint-iso}
把 \cref{eq:leg-contact-preint} 与 \cref{ch:imu} 的预积分 IMU 因子（\cref{eq:imu-preint-model,eq:imu-residual}）并排看，会发现它们\textbf{同构}：
\begin{center}\small
\begin{tabular}{lll}
\toprule
 & IMU 预积分（\cref{ch:imu}） & 接触/速度预积分（本章） \\
\midrule
高频测量 & 陀螺 $\tilde{\boldsymbol{\omega}}$、加计 $\tilde{\mathbf{a}}$ & 接触态 / 关节速度 $\dot{\mathbf{q}}$ \\
相对增量 & $\Delta\tilde{\mathbf{R}}_{ij},\Delta\tilde{\mathbf{v}}_{ij},\Delta\tilde{\mathbf{p}}_{ij}$ & $\Delta\tilde{\mathbf{B}}_{ij},\Delta\tilde{\mathbf{c}}_{ij}$ / $\Delta\tilde{\mathbf{t}}_{ij}$ \\
与端状态无关 & 扣重力/初速、左乘 $\mathbf{R}_i^\top$ & 接触系静止、左乘 $\mathbf{B}_i^\top$ / $\mathbf{R}_i^\top$ \\
噪声来源 & IMU 白噪声 $\boldsymbol{\eta}^g,\boldsymbol{\eta}^a$ & 打滑噪声 $\delta\boldsymbol{\theta},\delta\mathbf{d}$ / 编码器 $\boldsymbol{\eta}_q$ \\
协方差 & 一阶递推 $\boldsymbol{\Sigma}_{ij}$（\cref{eq:cov-recursion}） & 时间积分 $\boldsymbol{\Sigma}_w,\boldsymbol{\Sigma}_v$ / $A\boldsymbol{\Sigma}_v A^\top$ \\
进因子图 & 残差 $\mathbf{r}_{\mathcal{I}_{ij}}$，权 $\boldsymbol{\Sigma}_{ij}^{-1}$ & 残差 $\mathbf{r}_{\mathcal{C}}/\mathbf{r}_{\mathcal{V}}$，权 $\boldsymbol{\Lambda}=\boldsymbol{\Sigma}^{-1}$ \\
\bottomrule
\end{tabular}
\end{center}
这正是“融合而非拼贴”的体现：腿式预积分不是另起一套机理，而是把 \cref{ch:imu} 那套“高频测量 $\to$ 与初值无关的相对增量 $\to$ 因子图 $\to$ 信息矩阵加权”\emph{原封不动}地换一种传感器。读者只要握住 \cref{ch:imu} 的预积分，腿式因子就是“换了观测的同一道菜”。
\end{insight}

\begin{pitfall}[概念误区：接触预积分因子可以跨腿、跨站立相连用]
新手想法：“既然是相对增量，把任意两时刻接连起来不就行了？”\textbf{不行}。接触预积分因子\textbf{只对同一只站立腿、同一段站立相}有效——它的整个推导建在“这只脚这次踩着没动”上（\cref{eq:leg-contact-rot,eq:leg-contact-trans}）。脚一抬起，接触约束就断了，跨摆动相硬连会把腾空轨迹错算进去。现象：步态切换处轨迹突跳。正确做法：每段站立相内部用接触预积分，站立相之间靠 IMU 预积分（\cref{fig:leg-factor-graph} 的橙色因子）桥接；要在不同腿之间预积分，须显式建模\emph{接触切换}\cite{hartley2018legged}——这正是 Hartley 等后续工作处理“接触系切换动力学”要解决的。
\end{pitfall}

\subsection{速度预积分因子（四足）}
\label{subsec:leg-vel-preint}

\paragraph{动机：点足约束不了 $6$ 自由度，那就预积分线速度。}
点足四足无法用正运动学/接触预积分约束相对 $6$ 自由度（\cref{ins:leg-two-measurements}）。改用 \cref{eq:leg-vel-meas} 的线速度测量，把它\textbf{预积分}成约束机身相对位置变化的因子\cite{wisth2023vilens}。

\begin{derivation}[速度预积分残差]
设瞬时机身线速度由腿式 \cref{eq:leg-vel-meas} 得到、受高斯噪声 $\boldsymbol{\eta}_v$ 污染：
\begin{equation}
  \tilde{\mathbf{v}}=-\mathbf{J}_v(\mathbf{q})\dot{\mathbf{q}}-\boldsymbol{\omega}\times\mathbf{f}_p(\mathbf{q})+\boldsymbol{\eta}_v.
  \label{eq:leg-vel-noisy}
\end{equation}
设两时刻 $t_i,t_j$ 间机身线速度近似常值，对速度测量预积分（求和近似积分），得相对位置增量
\begin{equation}
  \Delta\tilde{\mathbf{t}}_{ij}=\Delta\mathbf{t}_{ij}+\delta\mathbf{p}_{ij}=\sum_{k=i}^{j-1}\big[\Delta\tilde{\mathbf{R}}_{ik}\,\tilde{\mathbf{v}}_k\,\Delta t\big]+\delta\mathbf{p}_{ij},
  \label{eq:leg-vel-sum}
\end{equation}
其中 $\delta\mathbf{p}_{ij}$ 是预积分速度噪声（与接触预积分的噪声地位相同）。于是\textbf{速度预积分因子}的残差与协方差：
\begin{align}
  \mathbf{r}_{\mathcal{V}}&=\mathbf{R}_i^\top(\mathbf{t}_j-\mathbf{t}_i)-\Delta\mathbf{t}_{ij},\label{eq:leg-vel-residual}\\
  \boldsymbol{\Sigma}_{\mathcal{V},ij}&=\sum_{k=i}^{j-1}\boldsymbol{\Sigma}_{\mathcal{V},ik}+\mathbf{A}\,\boldsymbol{\Sigma}_v\,\mathbf{A}^\top,\qquad \mathbf{A}=\Delta\tilde{\mathbf{R}}_{ik}\,\Delta t.\label{eq:leg-vel-cov}
\end{align}
\end{derivation}
残差 $\mathbf{r}_{\mathcal{V}}$ 把“状态给出的相对位移 $\mathbf{R}_i^\top(\mathbf{t}_j-\mathbf{t}_i)$”与“腿式速度积分出的位移 $\Delta\mathbf{t}_{ij}$”相减——与 \cref{ch:imu} 预积分位置残差 $\mathbf{r}_{\Delta\mathbf{p}_{ij}}$（\cref{eq:imu-residual} 第三行）\emph{结构完全一致}，只是这里的“速度”来自腿而非加计积分。协方差 \cref{eq:leg-vel-cov} 的递推形式也与 \cref{eq:cov-recursion} 同源（逐测量累加 $\mathbf{A}\boldsymbol{\Sigma}\mathbf{A}^\top$）。

\subsection{多腿同时站立怎么办}
\label{subsec:leg-multi}

前面每条因子只用一只腿。多腿同时站立带来冗余与鲁棒，但也可能\emph{冲突}（不同腿给矛盾信息）。三种处理：
\begin{itemize}
  \item \textbf{择优}：只选最可信的一只腿，弃其余\cite{carlone2026handbook}——最简单（SubT 中常见）。
  \item \textbf{独立处理}：把每条腿当独立测量各成一因子——接触位姿显式在状态里时最自然（\cref{eq:leg-state-quad}）。
  \item \textbf{平均}：把同时站立诸腿的速度测量平均成一条，降因子图/滤波的计算负担\cite{wisth2023vilens}。
\end{itemize}

\begin{insight}[滤波 vs 平滑：足式估计的频率分工]\label{ins:leg-filter-vs-smooth}
\cref{ch:eskf} 的滤波 vs 平滑对照（\cref{tab:filter-vs-opt}）在足式上有清晰的频率分工：\textbf{控制回路要高频}（约 $1\,\mathrm{kHz}$，防摔），故惯性 + 腿式融合多用滤波（EKF/UKF/InEKF）直接喂控制器；\textbf{建图要全局一致}，故慢速外感受（激光/相机）+ 腿式多用因子图平滑。卡尔曼滤波的一个固有局限是它需要过程模型，没有时只能拿常速度模型或 IMU 传播顶替；因子图把过程模型与测量模型\emph{一视同仁}地当作“状态与测量的关系”，更一般。二者按窗口长短连续过渡：只保两个相邻状态的因子图就退化成滤波（如 Two-State Implicit Filter, TSIF\cite{bloesch2013state}）；窗口加长则能回头修正历史，代价是计算。腿式预积分（本节）正是“拉长窗口又不爆变量数”的关键——和 IMU 预积分在 \cref{ch:imu} 扮演的角色一模一样。
\end{insight}

% ════════════════════════════════════════════════════════════════
\section{接触辅助不变卡尔曼滤波}
\label{sec:leg-inekf}

我们在 \cref{paper:inekf} 末尾许诺过：不变卡尔曼滤波（InEKF）的一个标志性应用，是足式机器人的接触辅助惯导。本节兑现它——这也是把腿式里程计接回全书理论主干最干净的一处：所需的群结构（$\mathrm{SE}_2(3)$，\cref{sec:se23}）、不变误差与群仿射一致性（\cref{paper:inekf}）全已在书中，本节只需把“接触约束”这一新观测织进去。

\paragraph{动机：为什么足式惯导特别适合 InEKF。}
回忆 \cref{paper:inekf} 的诊断：标准 EKF/ESKF 不一致的根源是“线性化雅可比依赖估计状态”，使误差与真值的耦合随轨迹漂移、虚假地制造可观测性。Barrau--Bonnabel 的解法\cite{barrau2017invariant} 是：若状态活在李群上、其确定性动力学\emph{群仿射}，则用\emph{不变误差}时，误差动力学\textbf{精确}自治（系统矩阵不含状态），线性化传播随之与轨迹无关（log-linear 性质）。惯导运动学（姿态由陀螺、速度由加计 + 重力、位置由速度）在 $\mathrm{SE}_2(3)$ 上恰是群仿射（\cref{sec:se23} 已说明速度必须装进群、不能当普通 $\mathbb{R}^3$ 拼在旁边，否则丢群仿射结构）。足式机器人正是“惯导 + 一堆接触约束”的场景，InEKF 的一致性优势在此尤为可贵——这也是 Hartley 等\cite{hartley2020contact} 当年提出接触辅助 InEKF 的初衷。

\subsection{把接触点装进状态群}
\label{subsec:leg-inekf-state}

\paragraph{增广 $\mathrm{SE}_2(3)$。}
接触辅助 InEKF 的状态在 \cref{sec:se23} 的扩展位姿群基础上，再为每只站立腿\textbf{增广}一列接触点位置 $\mathbf{c}_i\in\mathbb{R}^3$。沿用本书 $\mathrm{SE}_2(3)$ 的矩阵形式 $[\mathbf{R},\mathbf{r},\mathbf{v}]$（\cref{def:se23}：第 2 列位置 $\mathbf{r}$、第 3 列速度 $\mathbf{v}$），把接触点作为额外平移类列并入，得一个更大的矩阵李群元素（以单接触点为例）：
\begin{equation}
  \mathbf{X}=\begin{bmatrix}\mathbf{R} & \mathbf{r} & \mathbf{v} & \mathbf{c}\\\mathbf{0}^\top & 1 & 0 & 0\\\mathbf{0}^\top & 0 & 1 & 0\\\mathbf{0}^\top & 0 & 0 & 1\end{bmatrix},\qquad \mathbf{R}\in\SO(3),\ \mathbf{r},\mathbf{v},\mathbf{c}\in\mathbb{R}^3,
  \label{eq:leg-inekf-state}
\end{equation}
其中 $\mathbf{r}=\mathbf{t}_{wb}$ 为机身位置、$\mathbf{v}$ 为机身速度、$\mathbf{c}=\mathbf{c}_i$ 为接触点世界位置（多接触腿则并排多列，群相应增大，形如 $\mathrm{SE}_{K+2}(3)$）。这\textbf{正是 \cref{sec:se23} “$\mathrm{SE}_2(3)$ 几乎就是两份 $\SE(3)$”那条洞见（\cref{eq:se23-exp} 旁）的直接推广}：同一个旋转 $\mathbf{R}$ 搬运多个平移类量（位置、速度、各接触点），群的指数、伴随、雅可比的“形状”原样可知，不必从零推导。IMU 零偏 $\mathbf{b}^a,\mathbf{b}^g$ 仍以普通向量附在群旁——这会破坏严格群仿射，对应工程上的\emph{不完美 InEKF}（imperfect InEKF），与 \cref{paper:inekf} 的“实务限制”一段、\cref{sec:se23} 关于偏置的提醒完全一致。

\paragraph{右扰动与右不变误差。}
沿用全书右扰动 $\mathbf{X}\,\mathrm{Exp}(\delta\boldsymbol{\xi})$。误差取\emph{右不变} $\boldsymbol{\eta}=\hat{\mathbf{X}}\mathbf{X}^{-1}$（或左不变 $\mathbf{X}^{-1}\hat{\mathbf{X}}$，见 \cref{sec:lie-symmetry}）。\cref{paper:inekf} 的结论在此原样适用：群仿射 + 不变误差 $\Rightarrow$ 误差动力学自治、线性化传播与状态无关。

\subsection{接触约束作为不变观测}
\label{subsec:leg-inekf-meas}

\paragraph{动机：脚不动 = 一条天然的不变观测。}
站立腿的零速度约束（\cref{eq:leg-zero-vel}）与正运动学（\cref{eq:leg-fk-pos-rot}）给出观测：机身位姿 + 接触点位置应满足 $\mathbf{f}_p(\mathbf{q})=\mathbf{R}^\top(\mathbf{c}-\mathbf{r})$（脚在机身系的位置由编码器测得）。关键事实（Hartley 等\cite{hartley2020contact} 的核心）：在 \cref{eq:leg-inekf-state} 的增广群上，这条运动学观测恰好具有 InEKF 所要的\textbf{不变观测}形式——可写成增广状态作用在一个固定齐次向量上。于是它与惯导的群仿射动力学相容，InEKF 的一致性保证\emph{延续到带接触约束的足式惯导}：接触观测更新的线性化雅可比与估计状态无关，不破坏“误差传播与轨迹无关”这一性质。

\begin{derivation}[接触观测为何保持群结构（要点）]
把编码器测得的脚相对机身位置记作伪观测 $\tilde{\mathbf{y}}=\mathbf{f}_p(\tilde{\mathbf{q}})$。它与状态的关系 $\tilde{\mathbf{y}}=\mathbf{R}^\top(\mathbf{c}-\mathbf{r})+$ 噪声，可写成对增广状态 $\mathbf{X}$（\cref{eq:leg-inekf-state}）作用在一个\emph{固定齐次向量}上的形式——这正是 InEKF 理论里“右不变观测”的判据：观测 $\mathbf{Y}=\mathbf{X}^{-1}\mathbf{b}+$ 噪声（$\mathbf{b}$ 为常向量）。代入 \cref{eq:leg-inekf-state} 的 $\mathbf{X}^{-1}$（其旋转块 $\mathbf{R}^\top$、各平移列形如 $-\mathbf{R}^\top(\cdot)$，结构同 \cref{eq:se23-inverse}），第 $(\mathbf{c}-\mathbf{r})$ 组合恰好落在 $\mathbf{R}^\top(\mathbf{c}-\mathbf{r})$ 上。因此接触观测的更新雅可比\emph{不依赖估计状态}，与群仿射动力学一致——这是 InEKF 一致性在足式上成立的代数根由。完整滤波推导（增益、协方差更新）与 \cref{ch:eskf} 的 InEKF 同套，不再重述。
\end{derivation}

\begin{insight}[三线汇流：InEKF + $\mathrm{SE}_2(3)$ + 接触约束]\label{ins:leg-three-lines}
本节是全书三条伏笔的兑现处，值得停下来看清这条“一个作者通读后写的连贯论述”：
\begin{enumerate}
  \item \textbf{InEKF 线}（\cref{paper:inekf}）：当年讲 InEKF 时就预告了“足式接触辅助惯导”是其去向，本节把它落地——不必重新发明一致性理论，直接套用群仿射 + 不变误差。
  \item \textbf{$\mathrm{SE}_2(3)$ 线}（\cref{sec:se23}）：当年定义扩展位姿群时强调“识别生成元结构、闭式形状立即可知”，本节增广接触点 \cref{eq:leg-inekf-state} 正是这句话的直接应验——多一只脚就多一列平移类量，群机器原样迁移。
  \item \textbf{接触约束线}（本章 \cref{sec:leg-motion,sec:leg-contact}）：零速度 + 正运动学这条腿式观测，恰好是 InEKF 所要的右不变观测。
\end{enumerate}
三条线在足式平台上交于一点：腿式里程计不是孤立专题，而是“惯导一致性滤波”在多了一类本体感受约束时的自然形态。这正是把它放在 P2 末尾、紧接 \cref{ch:vio}/\cref{ch:lidar_slam} 之后的理由——同一套滤波/优化语言，再迁移一种平台。
\end{insight}

% ════════════════════════════════════════════════════════════════
\section{与外感受传感器融合}
\label{sec:leg-fusion}

腿式里程计提供一路额外的高频增量运动测量；它的主要用途是\textbf{改善里程计}，让其上的 SLAM 系统受益于节点间的低漂移边（回环时校正更平滑、不致大跳）。融合其他传感器（激光、相机）对滤波与平滑都很自然——前者多加测量、后者多加因子。但要把多路\emph{异质}测量（不同频率、噪声、失效率）融好并不平凡。本节复用本书已建立的耦合语言。

\paragraph{松耦合 vs 紧耦合：直接复用 VIO 的定义。}
\cref{ch:vio} 的 \cref{def:vio-coupling} 已为视觉惯性融合区分了松/紧耦合，这套语言原样适用于腿式融合：
\begin{itemize}
  \item \textbf{松耦合}：各子系统（视觉-惯性、腿-惯性、激光-惯性）并行独立运行，再对它们的\emph{输出}择优（triage）——若某模态破坏零均值高斯假设或彻底失效，就弃其输出选别的。这一思路由 DARPA SubT 大量催生\cite{carlone2026handbook}：在极端、传感器频繁退化的地下环境，运行多个互补子系统、随时挑当前最优，比孤注一掷于一个紧耦合系统更鲁棒。
  \item \textbf{紧耦合}：把各模态的\emph{原始}测量放进同一估计过程。Wisth 等\cite{wisth2023vilens} 用一个固定滞后平滑器（fixed-lag smoother）把腿式里程计、IMU、相机、激光放进\emph{同一}因子图；为应对模态间不一致或失效，\emph{择优直接发生在因子层面}——某模态失效就\emph{不往图里加}对应因子；为处理非零均值高斯噪声，因子里用\emph{鲁棒核函数}（\cref{ch:nlopt} 的 M-估计）。
\end{itemize}
\cref{def:vio-coupling} 旁那条洞见——“松耦合自举鲁棒、紧耦合精度高，是流水线上不同环节而非二选一”——在足式上同样成立。

\begin{insight}[腿-惯-视-激因子图：LIO-SAM/VIO 的自然推广]\label{ins:leg-multimodal}
\cref{ch:lidar_slam} 的 LIO-SAM 把 IMU 预积分、激光里程计、GPS、回环写成因子图四类因子；\cref{ch:vio} 的 VINS-Mono 把预积分与重投影因子放进同一滑窗。\textbf{腿式融合不过是在这张因子图上再加一类“腿式因子”}（接触预积分 \cref{eq:leg-contact-preint} 或速度预积分 \cref{eq:leg-vel-residual}），如 \cref{fig:leg-factor-graph} 所示。所有已学机制——预积分压缩高频测量、Schur 补边缘化老帧、FEJ 一致性、鲁棒核抗外点——\emph{原样可用}。这就是本章反复强调的“同一套数学换一种传感器”：腿-惯-视-激多模态因子图，是 LIO-SAM 因子图的自然推广，不是新发明。读者若已掌握 \cref{ch:vio}/\cref{ch:lidar_slam}，加一类腿式因子只是“多写一个残差与协方差”。
\end{insight}

% ════════════════════════════════════════════════════════════════
\section{开放问题与前沿}
\label{sec:leg-open}

前面的腿式里程计推导都默认了几条理想假设——刚体、刚性接触、不打滑——它们在实践中常被打破。本节先说这些假设失效的后果，再综述近年应对的趋势（多为学习驱动）。这一节面向研究，可选读。

\subsection{刚性假设的失效}
\label{subsec:leg-rigid-fail}

\paragraph{腿变形。}
正运动学 \cref{eq:leg-fk} 假设腿是刚体，算的是末端\emph{理想}位置。腿在负载下弯曲时，关节角虽变，真实末端位置却与理想不符，腿式测量出现\emph{偏差}（\cref{fig:leg-deform} 左：脚没动但腿弯了，正运动学误判为机身上抬）。腿越长问题越重，故人形比四足更受困扰。应对：(i) 建模腿载荷与柔度的关联（精确但复杂、难泛化）；(ii) 短时变形可靠分析力的剖面、在冲击段\emph{丢弃}测量\cite{bloesch2013state}；(iii) 最常用——在连杆上加 IMU、用其估连杆朝向以补偿\cite{carlone2026handbook}。

\begin{figure}[ht]
\centering
\begin{tikzpicture}[>=Stealth, font=\scriptsize]
  \draw[gray!50] (-0.4,-1.5) -- (2.0,-1.5);
  \draw[gray!50] (3.0,-1.5) -- (5.4,-1.5);
  % left: real (bent leg)
  \node[font=\small] at (0.8,0.9) {真实（腿弯）};
  \draw[fill=main!10, draw=main!60!black, rounded corners] (0.3,0.3) rectangle (1.3,0.6);
  \draw[thick, bend left=25] (0.8,0.3) to (0.8,-1.5);
  \fill[second!70!black] (0.8,-1.5) circle (1.4pt);
  % right: fk assumes rigid -> spurious up
  \node[font=\small] at (4.2,0.9) {正运动学（误判上抬）};
  \draw[fill=main!6, draw=main!50!black, rounded corners] (3.7,0.6) rectangle (4.7,0.9);
  \draw[thick] (4.2,0.6) -- (4.2,-1.5);
  \fill[second!70!black] (4.2,-1.5) circle (1.4pt);
  \draw[->, third!60!black, dashed] (4.2,0.45) -- (4.2,0.85) node[right]{虚假上抬};
\end{tikzpicture}
\caption{腿变形导致的偏差（仿 \cite{carlone2026handbook} 图 12.8 重绘）。脚踩在原地不动，但腿在负载下弯曲使关节角变化；刚体正运动学把这一变化误读成机身向上的运动。地面沉陷（脚陷入软土）会造成类似的“虚假上抬”，但破坏的是零速度假设而非刚体假设。}
\label{fig:leg-deform}
\end{figure}

\paragraph{非刚性接触与打滑。}
若地面松软、可塌陷，脚会陷入地面（\cref{fig:leg-deform} 同理）：接触\emph{先}于地面变形被检出，而接触点被假设静止，于是脚下沉被正运动学误读为机身上抬。这里破坏的是\emph{接触点零速度}假设（\cref{eq:leg-zero-vel}）而非刚体——接触点随地面变形向下移动。打滑同理：脚虽判为站立，其对地的力却违反摩擦圆锥 \cref{eq:leg-friction-cone}/CoP \cref{eq:leg-cop}，接触点其实在动。应对：需外感受传感器（如相机）让接触点速度在非退化运动/构型下\emph{可观}\cite{wisth2023vilens}——可把接触速度作额外状态显式跟踪，或取脚位置的导数。这是腿式里程计\emph{必须}与外感受融合（\cref{sec:leg-fusion}）的深层原因之一：本体感受单独看不出“地在动”。

\subsection{近年趋势}
\label{subsec:leg-trends}

\paragraph{学习接触估计。}
鉴于打滑、软地、腿变形难以解析建模，用数据驱动法估接触渐成趋势：学一个“何时建立接触”的二值信号。有监督（学接触分类器\cite{bloesch2013state}）、无监督（聚类，用关节 + 惯性等本体感受信号\cite{carlone2026handbook}）、以及深度网络（对更广的结构化/非结构化环境泛化更好\cite{hartley2020contact}）。基于视觉的触觉传感器（借计算机视觉的进展）是另一有前景的方向。回到 \cref{ins:leg-contact-assoc}：把接触估计这个“数据关联”做稳，是腿式里程计精度的命门，故它是头号前沿。

\paragraph{端到端学习与强化学习。}
高频腿式里程计与本体感受状态估计本是为基于模型的闭环运动控制而生；但强化学习（RL）的进展挑战了它的必要性。RL 运动控制器表明：运动只需机身\emph{速度}与\emph{相对重力的朝向}，二者可由标准本体感受估计器显式提供，甚至从关节 + 惯性原始数据\emph{隐式}学到\cite{carlone2026handbook}。这看似在质疑状态估计的必要——但仔细看，那是因为这类 RL 控制器的训练目标是\emph{跟踪速度指令}，而速度与重力朝向由惯性数据完全可观（\cref{ch:imu}），故可在训练中隐式估出。然而要完成更复杂的\emph{导航}（如机器人跑酷、奔向相对起点的目标），就必须显式的里程计来跟踪朝目标的进度——这说明状态估计在复杂运动与导航任务里仍不可或缺。一个折中方向是把状态估计作为运动策略学习的一部分，显式估机身速度、足高、接触态，有望给腿式里程计或 SLAM 的里程计因子提供更准的输入。

\paragraph{人形复兴。}
人形机器人承载了机器人学的部分梦想——拟人之身以无缝胜任为人设计的环境与工具。DRC（\cref{sec:leg-why}）是这方向的主要努力之一；其后多年重心转向更易控、更鲁棒的四足（工业巡检），直到 $2021$ 年 AI 乐观情绪与四足成功又重燃对人形的商业兴趣。人形把本章话题推向尚未充分探索的方向：长期、精确、可靠的\emph{全身}估计；脚之外还有躯干、手臂、手的\emph{间歇接触}（搬运包裹时）；与人共处所需的\emph{柔顺}——这要求放松本章的刚性接触假设。这些都是腿式状态估计的活跃前沿。

% ════════════════════════════════════════════════════════════════
\section{一个数值例：四足平地直行的速度里程计}
\label{sec:leg-example}

\begin{example}[四足单腿速度里程计的代数走查]\label{ex:leg-velocity}
为把 \cref{eq:leg-vel-meas} 落到实处，考虑一个理想化的四足在平地匀速直行：机身不转（$\boldsymbol{\omega}_b^w=\mathbf{0}$、$\mathbf{R}_{wb}=\mathbf{I}$），某站立腿的脚在机身系位置 $\mathbf{f}_p(\mathbf{q})=[0.20,\ -0.15,\ -0.40]^\top\,\mathrm{m}$（前 $20\,\mathrm{cm}$、右 $15\,\mathrm{cm}$、下 $40\,\mathrm{cm}$）。某瞬间该腿线雅可比与关节角速度乘积为 $\mathbf{J}_v(\mathbf{q})\dot{\mathbf{q}}=[-0.50,\ 0,\ 0]^\top\,\mathrm{m/s}$（脚相对机身向后 $0.5\,\mathrm{m/s}$）。

\textbf{第一步：套零速度约束。} 由 \cref{eq:leg-vel-meas}，机身速度
\[
  \mathbf{v}_b^w=-\underbrace{\boldsymbol{\omega}_b^w\times\mathbf{f}_p(\mathbf{q})}_{=\,\mathbf{0}\,(\text{不转})}-\mathbf{R}_{wb}\mathbf{J}_v(\mathbf{q})\dot{\mathbf{q}}=-[-0.50,0,0]^\top=[0.50,\ 0,\ 0]^\top\,\mathrm{m/s}.
\]
机身以 $0.5\,\mathrm{m/s}$ 前行——与“脚相对机身向后 $0.5\,\mathrm{m/s}$、脚踩死不动”的直觉一致（\cref{fig:leg-relpose}）。

\textbf{第二步：机身旋转时杆臂项不可省。} 现让机身以 $\boldsymbol{\omega}_b^w=[0,0,1.0]^\top\,\mathrm{rad/s}$ 偏航（绕竖轴 $1\,\mathrm{rad/s}$），其余不变。杆臂项
\[
  \boldsymbol{\omega}_b^w\times\mathbf{f}_p=[0,0,1.0]^\top\times[0.20,-0.15,-0.40]^\top=[0.15,\ 0.20,\ 0]^\top\,\mathrm{m/s},
\]
故 $\mathbf{v}_b^w=-[0.15,0.20,0]^\top-[-0.50,0,0]^\top=[0.35,\ -0.20,\ 0]^\top\,\mathrm{m/s}$。仅 $1\,\mathrm{rad/s}$ 的偏航，就让机身速度估计在 $x$ 上变化 $0.15$、在 $y$ 上凭空多出 $0.20\,\mathrm{m/s}$——与脚速度同量级。\textbf{结论}：若按 \cref{ins:imu-leverarm} 的告诫漏掉杆臂项 $-\boldsymbol{\omega}\times\mathbf{f}_p$，转弯时机身速度会系统性出错。这正是 \cref{eq:leg-vel-meas} 必须保留该项的原因，也呼应了 \cref{ch:imu} 杆臂效应一节——脚就是一段“几十厘米的杆臂”。

\textbf{第三步：噪声量级。} 设编码器位置噪声使 $\mathbf{J}_v\dot{\mathbf{q}}$ 有约 $1\%$ 相对误差（$\sim0.005\,\mathrm{m/s}$），由 \cref{eq:leg-vel-noise} 经雅可比传到 $\tilde{\mathbf{v}}$，构成速度因子 \cref{eq:leg-vel-residual} 协方差 $\boldsymbol{\Sigma}_{\mathcal{V}}$ 的来源。把若干步这样的速度因子预积分（\cref{eq:leg-vel-sum}），就得到一条约束相对位移的腿式因子——与一条 IMU 预积分因子并列进 \cref{fig:leg-factor-graph} 的因子图。
\end{example}

% ════════════════════════════════════════════════════════════════
\section*{常见误解汇总}
\begin{table}[H]\centering\small
\caption{腿式里程计常见误解与正解}
\label{tab:leg-misconceptions}
\begin{tabular}{p{0.42\textwidth} p{0.52\textwidth}}
\toprule
误解 & 正解 \\
\midrule
腿式里程计也得看环境（要相机/激光）& 纯本体感受：只用关节 + IMU + 力/接触，不看环境；唯一隐含假设是“脚不打滑”（\cref{ins:leg-proprioceptive}）\\
落地的接触点 $\mathbf{c}_i$ 像视觉路标一样永久固定 & 只在当次站立相固定；脚再落地是新接触，须重置 $\mathbf{c}_i$ 并放大其不确定度 \\
接触/速度预积分是与 IMU 不同的新机理 & 与 IMU 预积分\emph{同构}：高频测量 $\to$ 与初值无关的相对增量 $\to$ 因子图 $\to$ $\boldsymbol{\Lambda}$ 加权（\cref{ins:leg-preint-iso}）\\
接触预积分因子可跨腿、跨站立相连用 & 只对同一站立腿、同一站立相有效；跨相靠 IMU 桥接，跨腿须显式建模接触切换 \\
点足四足也能用正运动学定 $6$ 自由度相对位姿 & 点足脚可绕接触点自转，无姿态信息；只能估线速度 $\Rightarrow$ 走速度预积分（\cref{ins:leg-two-measurements}）\\
机身转弯时速度估计不必管脚的位置 & 杆臂项 $-\boldsymbol{\omega}\times\mathbf{f}_p$ 与脚速度同量级，转弯漏掉它会系统性出错（\cref{ex:leg-velocity}）\\
InEKF 是 VIO 专用 & 足式接触辅助惯导是 InEKF 的标志性落地；接触点装进增广 $\mathrm{SE}_2(3)$ 即可（\cref{sec:leg-inekf}）\\
脚下沉/打滑导致的“鼓包”是腿弯了 & 腿变形破坏\emph{刚体}假设；地面沉陷/打滑破坏\emph{接触点零速度}假设——两者现象像、根因不同（\cref{subsec:leg-rigid-fail}）\\
固定一个 $f_z$ 阈值就能判接触 & 阈值随步态/地形/负载变；需滞回 + 力的时间模式，或学习的接触分类器 \\
\bottomrule
\end{tabular}
\end{table}

% ════════════════════════════════════════════════════════════════
\section*{本章小结}

\begin{table}[H]\centering\small
\caption{符号表}
\begin{tabular}{lll}
\toprule
符号 & 含义 & 首次出现 \\
\midrule
$\mathcal{F}^w,\mathcal{F}^b,\mathcal{F}^f,\mathcal{F}^k$ & 世界 / 机身 / 脚 / 接触系 & \cref{subsec:leg-frames} \\
$\mathbf{q},\dot{\mathbf{q}}$ & 关节角 / 关节角速度（编码器） & \cref{eq:leg-state-raw} \\
$\mathbf{fk}(\mathbf{q}),\mathbf{f}_p,\mathbf{f}_R$ & 正运动学 / 其位置、姿态分量 & \cref{eq:leg-fk,eq:leg-fk-pos-rot} \\
$\mathbf{J}(\mathbf{q}),\mathbf{J}_v,\mathbf{J}_\omega$ & 雅可比 / 其线、角部分 & \cref{eq:leg-jacobian} \\
$\mathbf{T}_{b'b}$ & 机身相对位姿测量 & \cref{eq:leg-fk-meas} \\
$\mathbf{c}_i,\mathbf{B}_i,\mathbf{C}_i$ & 接触点位置 / 姿态 / 位姿（世界系） & \cref{eq:leg-state-quad,eq:leg-Tc-expand} \\
$\mathbf{v}_b^w,\boldsymbol{\omega}_b^w$ & 机身线 / 角速度 & \cref{eq:leg-vel-meas} \\
$\mathbf{f}=[f_x,f_y,f_z]^\top,\mu$ & 地面反力 / 摩擦系数 & \cref{eq:leg-friction-cone} \\
$\mathrm{CoP}$ & 压力中心 & \cref{eq:leg-cop} \\
$\boldsymbol{\eta}_q,\boldsymbol{\Sigma}_q,\mathbf{J}_c$ & 编码器噪声 / 协方差 / 接触系雅可比 & \cref{eq:leg-encoder-noise,eq:leg-fp-noise} \\
$\mathbf{r}_{\mathcal{F}},\mathbf{r}_{\mathcal{C}},\mathbf{r}_{\mathcal{V}}$ & 正运动学 / 接触 / 速度预积分残差 & \cref{eq:leg-fk-residual,eq:leg-contact-preint,eq:leg-vel-residual} \\
$\delta\boldsymbol{\theta}_{ij},\delta\mathbf{d}_{ij},\delta\mathbf{p}_{ij}$ & 预积分接触/速度噪声 & \cref{eq:leg-contact-rot,eq:leg-vel-sum} \\
$\mathbf{X}\in\mathrm{SE}_{K+2}(3)$ & 接触辅助 InEKF 增广状态 & \cref{eq:leg-inekf-state} \\
\bottomrule
\end{tabular}
\end{table}

\begin{table}[H]\centering\small
\caption{定理/公式速查表}
\begin{tabular}{lll}
\toprule
公式 & 一句话说明 & 对应 \\
\midrule
正运动学 \cref{eq:leg-fk} & 关节角 $\to$ 脚位姿（同机械臂 FK） & \cref{subsec:leg-fk} \\
雅可比 \cref{eq:leg-jacobian} & 关节角速度 $\to$ 脚速度；稀疏块 & \cref{subsec:leg-fk} \\
相对位姿 \cref{eq:leg-fk-meas} & 接触静止 $\Rightarrow$ $\mathbf{fk}(\mathbf{q}')^{-1}\mathbf{fk}(\mathbf{q})$ & \cref{subsec:leg-relpose} \\
腿式观测 \cref{eq:leg-fk-pos-rot} & $\mathbf{f}_p=\mathbf{R}^\top(\mathbf{c}_i-\mathbf{t})$ 联系状态与测量 & \cref{subsec:leg-relpose} \\
机身速度 \cref{eq:leg-vel-meas} & 零速度约束 $\Rightarrow$ $-\boldsymbol{\omega}\times\mathbf{f}_p-\mathbf{R}\mathbf{J}_v\dot{\mathbf{q}}$ & \cref{subsec:leg-velocity} \\
摩擦圆锥 \cref{eq:leg-friction-cone} & 不打滑 $\Leftrightarrow$ 力落锥内；同控制侧 & \cref{subsec:leg-friction} \\
GRF 反演 \cref{eq:leg-grf-from-torque} & 关节力矩 + 动力学 $\to$ 地面反力 & \cref{subsec:leg-grf} \\
编码器传播 \cref{eq:leg-fp-noise} & 雅可比一阶近似保高斯 & \cref{subsec:leg-encoder-noise} \\
正运动学因子 \cref{eq:leg-fk-residual} & $\mathrm{Log}(\mathbf{C}_i^{-1}\mathbf{T}\,\mathbf{fk}(\tilde{\mathbf{q}}))$ & \cref{subsec:leg-fk-factor} \\
接触预积分 \cref{eq:leg-contact-preint} & 约束接触位姿变化；同 IMU 预积分构 & \cref{subsec:leg-contact-preint} \\
速度预积分 \cref{eq:leg-vel-residual} & 积线速度成位移约束（点足四足） & \cref{subsec:leg-vel-preint} \\
增广 $\mathrm{SE}_2(3)$ \cref{eq:leg-inekf-state} & 接触点并入群 $\Rightarrow$ 接触辅助 InEKF & \cref{subsec:leg-inekf-state} \\
\bottomrule
\end{tabular}
\end{table}

\begin{table}[H]\centering\small
\caption{知识点总表}
\begin{tabular}{clll}
\toprule
\# & 知识点 & 核心要点 & 难度 \\
\midrule
1 & 参考系与状态 & 世界/机身/脚/接触系；状态按 $\mathrm{SE}_2(3)$ 立、按需增广 & $\bigstar\bigstar$ \\
2 & 腿运动学 $\mathbf{fk}/\mathbf{J}$ & 同机械臂 FK；雅可比稀疏块；与控制共享 & $\bigstar\bigstar$ \\
3 & 相对位姿估计 & 接触静止几何；增广接触点绕开三足约束 & $\bigstar\bigstar\bigstar$ \\
4 & 速度估计 & 接触点零速度约束 + 杆臂项 & $\bigstar\bigstar\bigstar$ \\
5 & 接触估计 & 摩擦锥/CoP、阈值化 $f_z$、力矩反演 GRF & $\bigstar\bigstar\bigstar$ \\
6 & 编码器噪声传播 & 雅可比一阶近似保高斯 & $\bigstar\bigstar$ \\
7 & 接触/速度预积分因子 & IMU 预积分的同构平行实例 & $\bigstar\bigstar\bigstar\bigstar$ \\
8 & 接触辅助 InEKF & 接触点入增广 $\mathrm{SE}_2(3)$；右不变观测保一致性 & $\bigstar\bigstar\bigstar\bigstar$ \\
9 & 多传感融合 & 松/紧耦合；腿-惯-视-激因子图 & $\bigstar\bigstar\bigstar$ \\
10 & 开放问题与前沿 & 腿变形/沉陷/打滑；学习接触、RL、人形 & $\bigstar\bigstar\bigstar$ \\
\bottomrule
\end{tabular}
\end{table}

% ════════════════════════════════════════════════════════════════
\section*{延伸阅读}
\begin{itemize}
  \item \textbf{系统综述}：SLAM Handbook\cite{carlone2026handbook} 第 12 章（Camurri \& Mattamala）——腿式里程计的参考系、运动学、接触估计、因子图与开放问题（本章框架与分类的出处）。
  \item \textbf{滤波奠基}：Bloesch 等\cite{bloesch2013state}——基于 EKF 的全地形四足运动学-惯性状态估计，奠定本体感受滤波框架（接触点增广、零速度约束）。
  \item \textbf{接触辅助不变滤波}：Hartley 等\cite{hartley2020contact}（IJRR 2020）——接触辅助 InEKF、不完美 InEKF，把 \cref{paper:inekf} 的不变滤波落到足式惯导（本章 \cref{sec:leg-inekf} 的出处）。
  \item \textbf{接触预积分}：Hartley 等\cite{hartley2018legged}——人形的前向运动学因子与接触预积分因子，并处理接触系切换（本章 \cref{subsec:leg-fk-factor,subsec:leg-contact-preint}）。
  \item \textbf{速度预积分与紧耦合}：Wisth 等\cite{wisth2023vilens}（VILENS，T-RO 2023）——四足的速度预积分因子，把腿-惯-视-激放进同一固定滞后平滑器（本章 \cref{subsec:leg-vel-preint,sec:leg-fusion}）。
  \item \textbf{扩展位姿与不变滤波地基}：Barrau--Bonnabel\cite{barrau2017invariant}（InEKF 与群仿射/log-linear 性质）；本书 \cref{sec:se23}（$\mathrm{SE}_2(3)$ 权威定义）与 \cref{paper:inekf}（InEKF 推导）。
  \item \textbf{预积分理论}：Forster 等\cite{forster2017preintegration}（流形 IMU 预积分，本章腿式预积分的母版）；本书 \cref{ch:imu}。
  \item \textbf{机器人运动学/动力学}：Lynch \& Park《Modern Robotics》\cite{lynch2017modern}（旋量积正运动学、雅可比）；Featherstone《Rigid Body Dynamics Algorithms》\cite{featherstone2008rigid}（递推 Newton--Euler、浮基动力学）；本书 \cref{sec:ctrl-legged} 足式控制。
\end{itemize}

\section*{本章与其它章节的关系}
\begin{table}[H]\centering\small
\begin{tabular}{lll}
\toprule
章节 & 关系 & 本章承接/铺垫 \\
\midrule
\cref{paper:inekf} InEKF & 兑现其“足式接触辅助惯导”前向许诺 & 群仿射 + 不变误差 $\Rightarrow$ 一致性 \\
\cref{sec:se23} $\mathrm{SE}_2(3)$ & 接触点增广进群（\cref{eq:leg-inekf-state}） & “多一列平移类量”的洞见直接推广 \\
\cref{ch:imu} IMU 预积分 & 接触/速度预积分是其同构平行实例 & 高频测量 $\to$ 相对增量 $\to$ 因子图 \\
\cref{ch:vio} VIO & 复用松/紧耦合定义（\cref{def:vio-coupling}） & 腿式因子加进同套因子图 \\
\cref{ch:lidar_slam} 激光 SLAM & 腿-惯-视-激是 LIO-SAM 因子图的推广 & 紧耦合因子图、鲁棒核 \\
\cref{sec:ctrl-legged} 足式控制 & 共享 $\mathbf{fk}/\mathbf{J}$ 与摩擦锥（互为表里） & 估计给控制提供速度/朝向 \\
\cref{ch:radar} 雷达 / \cref{ch:event} 事件 & 同属“新型传感器与平台”章群 & 同一套数学迁移新平台 \\
\bottomrule
\end{tabular}
\end{table}

\section*{故障排查手册}
\begin{table}[H]\centering\small
\begin{tabular}{p{0.22\textwidth} p{0.30\textwidth} p{0.38\textwidth}}
\toprule
症状 & 可能原因 & 排查步骤 \\
\midrule
脚抬起瞬间位姿突跳 & 接触切换时未重置 $\mathbf{c}_i$ & 1.确认摆动转站立时用机身位姿 + FK 重置 $\mathbf{c}_i$；2.重置时放大其协方差；3.核对接触检测时序（\cref{subsec:leg-relpose}）\\
$z$ 轴系统性“鼓包”/漂高 & 腿变形或地面沉陷/打滑 & 1.区分刚体失效 vs 零速度失效（\cref{subsec:leg-rigid-fail}）；2.冲击段丢测量或加连杆 IMU；3.软地上引入外感受让接触速度可观 \\
转弯时速度估计明显偏 & 漏杆臂项 $-\boldsymbol{\omega}\times\mathbf{f}_p$ & 1.核对 \cref{eq:leg-vel-meas} 含杆臂项；2.用 \cref{ex:leg-velocity} 第二步量级自检；3.确认陀螺 $\boldsymbol{\omega}$ 与脚位置 $\mathbf{f}_p$ 同一时刻 \\
站立/摆动误判频繁 & 固定 $f_z$ 阈值不适配步态/地形 & 1.改滞回双阈值；2.用力的时间模式而非瞬时值；3.考虑学习接触分类器（\cref{subsec:leg-fz}）\\
接触预积分因子致图发散 & 跨站立相/跨腿误连 & 1.确认因子只在同一站立腿同一相内；2.相间用 IMU 因子桥接；3.接触切换显式建模（\cref{subsec:leg-contact-preint}）\\
紧耦合某模态失效拖垮全图 & 未做因子级择优 / 无鲁棒核 & 1.失效模态\emph{不加}对应因子；2.因子用鲁棒核；3.必要时退松耦合择优（\cref{sec:leg-fusion}）\\
InEKF 一致性不如预期 & 偏置增广破坏群仿射 & 1.确认用不完美 InEKF 近似处理偏置；2.核对接触观测写成右不变形式；3.对照 \cref{paper:inekf} 的实务限制 \\
\bottomrule
\end{tabular}
\end{table}

\section*{API 速查表}
\begin{table}[H]\centering\small
\begin{tabular}{p{0.30\textwidth} p{0.62\textwidth}}
\toprule
功能 & 接口 / 要点（一句话） \\
\midrule
正运动学 / 雅可比 & Pinocchio \texttt{forwardKinematics}、\texttt{computeJointJacobians}；或 RBDL、KDL（给 $\mathbf{fk},\mathbf{J}$） \\
浮基动力学 / GRF & Pinocchio \texttt{rnea}/\texttt{aba}（递推 Newton--Euler / ABA）；\cref{eq:leg-grf-from-torque} 反演 \\
李群 / 右扰动 & Sophus \texttt{SO3d/SE3d::exp/log}；接触/速度因子残差用 $\mathrm{Log}$ + 右扰动 \\
因子图 / 平滑 & GTSAM \texttt{ISAM2}/\texttt{FixedLagSmoother}、\texttt{ImuFactor}（预积分）、自定义 leg factor（\texttt{NoiseModelFactor}） \\
接触辅助 InEKF & 开源 \texttt{invariant-ekf}（Hartley 等）；状态在增广 $\mathrm{SE}_K(3)$ 上 \\
鲁棒核 & GTSAM \texttt{noiseModel::Robust}（Huber/Tukey）；紧耦合多模态抗外点 \\
开源系统 & Cheetah-Software（MIT）、Pronto/VILENS（牛津）、Cerberus、ANYmal 状态估计栈 \\
\bottomrule
\end{tabular}
\end{table}

\section*{研究实践建议}
\begin{itemize}
  \item \textbf{初学者}：用 Pinocchio 对一个开源四足 URDF 算单腿 $\mathbf{fk}/\mathbf{J}$，在平地匀速数据上验证 \cref{eq:leg-vel-meas} 给出的机身速度，并复算 \cref{ex:leg-velocity}。
  \item \textbf{进阶}：实现一条速度预积分因子（\cref{eq:leg-vel-residual,eq:leg-vel-cov}），与一条 GTSAM \texttt{ImuFactor} 一起放进固定滞后平滑器，逐项核对其残差/协方差与 \cref{ch:imu} 的预积分同构（\cref{ins:leg-preint-iso}）。
  \item \textbf{有余力}：把接触点增广进 $\mathrm{SE}_2(3)$ 跑一个接触辅助 InEKF（\cref{sec:leg-inekf}），在打滑/软地序列上比较它与 ESKF 的一致性；再加一路相机让接触点速度可观，观察 $z$ 轴“鼓包”是否被抑制（\cref{subsec:leg-rigid-fail}）。
\end{itemize}

\section*{习题}
\begin{enumerate}
  \item （运动学）写出“接触系静止 $\Rightarrow$ $\mathbf{T}_{b'b}=\mathbf{fk}(\mathbf{q}')^{-1}\mathbf{fk}(\mathbf{q})$”（\cref{eq:leg-fk-meas}）的完整推导，指明哪一步用到“脚系与接触系在站立期重合”。再说明：为何点足四足靠它定 $6$ 自由度相对位姿需至少三足同时站立？
  \item （速度约束）从接触点零速度 \cref{eq:leg-zero-vel} 与刚体运动学出发，独立推出机身速度 \cref{eq:leg-vel-meas}，并指出杆臂项 $-\boldsymbol{\omega}\times\mathbf{f}_p$ 的物理来源。用 \cref{ex:leg-velocity} 的数据验证你的公式。
  \item （接触估计）解释为何“阈值化 $f_z$”能同时（近似）满足摩擦圆锥 \cref{eq:leg-friction-cone} 与 CoP 条件 \cref{eq:leg-cop}。再讨论固定阈值在快步着地与软地下坡两种情形各为何失效，给出改进。
  \item （预积分同构）把接触预积分残差 \cref{eq:leg-contact-preint} 与 \cref{ch:imu} 的 IMU 预积分残差 \cref{eq:imu-residual} 逐项对照，论证二者同构（填充 \cref{ins:leg-preint-iso} 表的每一行）。哪一项在腿式里没有对应、为什么？
  \item （只在同一站立相）论证接触预积分因子为何只对同一站立腿、同一站立相有效（提示：\cref{eq:leg-contact-rot,eq:leg-contact-trans} 的推导前提）。要在不同腿之间预积分需补什么模型？
  \item （InEKF 增广）仿 \cref{eq:se23-inverse} 写出增广状态 \cref{eq:leg-inekf-state} 的逆 $\mathbf{X}^{-1}$，并验证接触观测 $\mathbf{f}_p=\mathbf{R}^\top(\mathbf{c}-\mathbf{r})$ 可写成右不变形式 $\mathbf{X}^{-1}\mathbf{b}$（$\mathbf{b}$ 为常齐次向量）。由此说明其更新雅可比为何不依赖估计状态（呼应 \cref{paper:inekf}）。
  \item （编码器传播）由 \cref{eq:leg-vel-noise} 论证：速度测量同时受编码器位置噪声 $\boldsymbol{\eta}_q$ 与角速度噪声 $\boldsymbol{\eta}_{\dot{q}}$ 影响，且二者都可经已知雅可比线性传到 $\tilde{\mathbf{v}}$ 的协方差。这与 \cref{ch:imu} 把 IMU 白噪声经 $\mathbf{A},\mathbf{B}$ 传播（\cref{eq:cov-recursion}）在思想上有何共通？
  \item （融合与开放）某四足在覆雪松软坡地行走，纯腿-惯里程计的 $z$ 轴稳定漂高。诊断是腿变形还是地面沉陷？分别破坏了本章哪条假设？为消除该漂移，应加哪类外感受传感器、让什么量变得可观（\cref{subsec:leg-rigid-fail}）？
\end{enumerate}

% ════════════════════════════════════════════════════════════════
\section*{本章版本信息}
\begin{finenote}
本章系统吸收 SLAM Handbook 第 12 章（Camurri \& Mattamala, \emph{Leg Odometry for SLAM}），并融合 Bloesch 等的运动学-惯性滤波\cite{bloesch2013state}、Hartley 等的接触辅助 InEKF 与接触预积分\cite{hartley2020contact,hartley2018legged}、Wisth 等的速度预积分\cite{wisth2023vilens}，全章自包含、右扰动主线。腿运动学 $\mathbf{fk}/\mathbf{J}$、接触系静止几何、零速度约束、接触/速度预积分因子均独立重写，可不依赖源书读懂。本章是全书三条已有线的汇流点：开篇兑现 \cref{paper:inekf} 的“足式接触辅助惯导”前向许诺，接触辅助 InEKF 直接增广 \cref{sec:se23} 的 $\mathrm{SE}_2(3)$，腿式预积分写成 \cref{ch:imu} IMU 预积分的同构平行实例，并与 \cref{sec:ctrl-legged} 的足式控制共享 $\mathbf{fk}/\mathbf{J}$ 与摩擦锥互为表里。属 P2 末尾“新型传感器与平台里程计”章群，与雷达 SLAM（\cref{ch:radar}）、事件相机 SLAM（\cref{ch:event}）为兄弟章。教学层遵循 v5.0 算法工程教学规范。
\end{finenote}
